\documentclass[11pt,a4paper]{article}
\usepackage[utf8]{inputenc}
\usepackage[T1]{fontenc}
\usepackage[spanish,es-tabla,es-noshorthands]{babel}
\usepackage[top=2.3cm,bottom=2.3cm,left=2.2cm,right=2.2cm,headheight=14pt]{geometry}
\usepackage{amsmath,amssymb,mathtools}
\usepackage{icomma}
\spanishplainpercent
\usepackage{graphicx,booktabs,array,multirow,multicol,enumitem}
\usepackage[table]{xcolor}
\usepackage[most]{tcolorbox}
\usepackage{titlesec}
\usepackage{fancyhdr}
\usepackage{needspace}
\usepackage[hidelinks]{hyperref}

\definecolor{azul}{RGB}{35,85,140}
\definecolor{azulclaro}{RGB}{238,244,251}
\definecolor{rojo}{RGB}{180,40,40}
\definecolor{rojoclaro}{RGB}{252,240,238}
\definecolor{verde}{RGB}{30,110,50}
\definecolor{verdeclaro}{RGB}{238,248,240}
\definecolor{gris}{RGB}{245,245,245}

\setlength{\parindent}{0pt}
\setlength{\parskip}{4pt}
\setlist{itemsep=1pt,topsep=2pt}

\titleformat{\section}{\Large\bfseries\color{azul}}{Bloque \thesection\ ---}{0.4em}{}
\titleformat{\subsection}{\large\bfseries\color{azul}}{\thesubsection}{0.6em}{}
\titlespacing*{\section}{0pt}{10pt}{8pt}
\titlespacing*{\subsection}{0pt}{14pt}{6pt}

\pagestyle{fancy}
\fancyhf{}
\lhead{\small Probabilidad y Estadística -- Guía paso a paso del segundo parcial}
\rhead{\small Procedimientos generales}
\cfoot{\thepage}
\renewcommand{\headrulewidth}{0.4pt}

\newtcolorbox{cuando}{colback=verdeclaro,colframe=verde!70,boxrule=0.5pt,arc=2pt,left=6pt,right=6pt,top=3pt,bottom=3pt,
  fonttitle=\bfseries\small,coltitle=white,colbacktitle=verde!75,title={Cuándo se usa},breakable}
\newtcolorbox{pasos}{colback=white,colframe=azul!70,boxrule=0.7pt,arc=2pt,left=6pt,right=6pt,top=3pt,bottom=3pt,
  fonttitle=\bfseries\small,coltitle=white,colbacktitle=azul!85,title={Paso a paso},breakable}
\newtcolorbox{modelo}[1]{colback=gris,colframe=black!40,boxrule=0.5pt,arc=2pt,left=6pt,right=6pt,top=3pt,bottom=3pt,
  fontupper=\small,fonttitle=\bfseries\small,coltitle=black,colbacktitle=black!10,title={Modelo resuelto --- #1},breakable}
\newtcolorbox{ojo}[1][Errores típicos]{colback=rojoclaro,colframe=rojo!80,boxrule=0.5pt,arc=2pt,left=6pt,right=6pt,top=3pt,bottom=3pt,
  fonttitle=\bfseries\small,coltitle=white,colbacktitle=rojo!80,title={#1},breakable}
\newtcolorbox{hoja}[2][]{colback=white,colframe=azul!70,boxrule=0.6pt,arc=2pt,left=4pt,right=4pt,top=2pt,bottom=2pt,
  fonttitle=\bfseries\small,colbacktitle=azul!85,coltitle=white,title={#2},breakable,#1}

\newenvironment{pasolista}{\begin{enumerate}[label=\textbf{Paso \arabic*.},leftmargin=*,align=left,itemsep=3pt]}{\end{enumerate}}
\newcommand{\E}{\mathrm{E}}
\newcommand{\V}{\mathrm{V}}
\newcommand{\Cov}{\mathrm{Cov}}
\newcommand{\Nor}{\mathcal{N}}
\newcommand{\xb}{\bar{x}}
\newcommand{\Xb}{\bar{X}}
\newcommand{\chis}{\chi^2}

\begin{document}

\begin{center}
{\LARGE\bfseries\color{azul} Probabilidad y Estadística}\\[4pt]
{\large Guía de procedimientos paso a paso --- Segundo Parcial}\\[2pt]
{\normalsize Cómo resolver cada tipo de ejercicio, con modelos de parciales anteriores y hoja de fórmulas}\\[8pt]
{\normalsize Caetano Gimenez}\\[1pt]
{\small UTN - Facultad Regional Haedo \ $\cdot$\ Noviembre 2026}
\end{center}

\begin{tcolorbox}[colback=white,colframe=azul!60,boxrule=0.5pt,arc=2pt,title={\bfseries Cómo usar esta guía},colbacktitle=azul!80,fonttitle=\small]
\small
\begin{enumerate}[leftmargin=14pt]
\item Buscá en la tabla de abajo qué pide el enunciado y andá al procedimiento indicado.
\item Cada procedimiento tiene: \textbf{cuándo se usa} (verde), el \textbf{paso a paso general} (azul), un \textbf{modelo resuelto} con un ejercicio de los parciales y modelos de la cátedra (gris) y los \textbf{errores típicos} (rojo).
\item Los pasos son generales: sirven para cualquier ejercicio de ese tipo, cambiando los números.
\item Al final está la \textbf{hoja de fórmulas completa}.
\end{enumerate}
Notación: $\Phi(z)=P(Z\le z)$. Los valores críticos dejan el área \emph{a la derecha}: $P(Z>z_\alpha)=\alpha$, $P(T_\nu>t_{\alpha,\nu})=\alpha$, $P(\chis_\nu>\chis_{\alpha,\nu})=\alpha$. Como pide el parcial, conviene usar al menos 4 decimales.
\end{tcolorbox}

\subsection*{¿Qué tipo de ejercicio es?}
\begin{center}
\small
\renewcommand{\arraystretch}{1.18}
\begin{tabular}{@{}>{\raggedright\arraybackslash}p{10.6cm}l@{}}
\toprule
Si el enunciado pide o dice\ldots & Procedimiento\\\midrule
``Hallar la matriz / tabla / función de probabilidad conjunta'' & \ref{p:conjunta}\\
``Si se sabe que\ldots, ¿cuál es la probabilidad de\ldots?'', $P(A\,|\,B)$, condicional & \ref{p:condicional}\\
``¿Son independientes? Justificar'' & \ref{p:indep}\\
``Covarianza'', ``coeficiente de correlación'', $\E(XY)$ & \ref{p:cov}\\
``$f(x,y)=\dots$'' (densidad conjunta), integrales dobles & \ref{p:continuo}\\
``Varianza de $W=aX+bY+c$'', ``media y distribución de $w$'' & \ref{p:comblin}\\
``Total / suma / promedio de $n$\ldots'', ``consumo anual'', ``caja de 24'', ``¿cuánto debe valer $n$?'' & \ref{p:tcl}\\
``Por lo menos $k$ de $n$'', binomial con $n$ grande, ``resultado aproximado'' & \ref{p:aprox}\\
``Estimador insesgado'', ``¿cuál es preferible?'', ``varianza mínima'', ECM & \ref{p:estim}\\
``Estimar la media con confianza\ldots'' & \ref{p:icmu}\\
``Intervalo para la varianza / el desvío'' & \ref{p:icvar}\\
``Estimar la proporción'', ``¿con qué nivel de confianza se hizo el IC?'' & \ref{p:icp}\\
``¿Cuántos hay que muestrear para estimar\ldots?'', ``¿cuántos más?'' & \ref{p:nest}\\
``¿Hay evidencia\ldots?'', ``plantear la prueba de hipótesis'' & \ref{p:general} y \ref{p:pmu}--\ref{p:pp}\\
``Calcular el valor $P$'', ``¿a qué conclusión se llega si $P=\dots$?'' & \ref{p:valorp}\\
``Nivel de significación de la regla'', ``probabilidad de error tipo II'', ``probabilidad de detectar / adoptar'' & \ref{p:ab}\\
``Indicar el gráfico'', ``¿es una curva de potencia?'', ``deducir del gráfico'' & \ref{p:curvas}\\
``Tamaño de muestra para que la prueba detecte\ldots con probabilidad\ldots'', ``¿alcanza con $n$?'' & \ref{p:nprueba}\\
``Comparar el intervalo con la conclusión de la prueba'' & \ref{p:dual}\\
``Verdadero o falso'', ``demostrar'', ``deducir'', ``enunciar'' & \ref{p:vf}--\ref{p:interp}\\
\bottomrule
\end{tabular}
\end{center}

\newpage
{\small\setcounter{tocdepth}{2}
\begin{multicols}{2}
\tableofcontents
\end{multicols}}

% =====================================================================
\newpage
\section{Variables aleatorias bidimensionales}
% =====================================================================

\subsection{Armar la distribución conjunta (caso discreto)}\label{p:conjunta}
\begin{cuando}
El enunciado describe un experimento (extracciones de una caja, nodos observados de a uno, dados, monedas, clientes que eligen cajas, componentes que funcionan) y pide la \textbf{matriz}, \textbf{tabla} o \textbf{función de probabilidad conjunta} $p(x,y)$.
\end{cuando}
\begin{pasos}
\begin{pasolista}
\item \textbf{Definir las variables.} Escribí qué cuenta $X$ y qué cuenta $Y$, y sus \textbf{valores posibles} (rango). Ejemplo: ``$X$: cantidad de azules entre las 2 extraídas, $X\in\{0,1,2\}$''.
\item \textbf{Reconocer el mecanismo del experimento}, porque define cómo se calculan las probabilidades:
\begin{itemize}
\item \emph{Sin reposición y el orden no importa} (se sacan $n$ juntos): \textbf{conteo combinatorio},
$p(x,y)=\dfrac{\binom{N_1}{x}\binom{N_2}{y}\binom{N_3}{n-x-y}}{\binom{N}{n}}$.
\item \emph{De a uno, con regla de parada o cuando el orden importa}: \textbf{árbol de secuencias} y \textbf{regla del producto}, $P(A_1A_2A_3)=P(A_1)\,P(A_2|A_1)\,P(A_3|A_1A_2)$.
\item \emph{Ensayos independientes} (dados, monedas, componentes, clientes): producto de las probabilidades de cada ensayo.
\item \emph{El enunciado dice que $X$ e $Y$ son independientes}: directamente $p(x,y)=p_X(x)\,p_Y(y)$.
\end{itemize}
\item \textbf{Listar todos los resultados posibles} (secuencias o combinaciones). Para cada uno calculá su probabilidad y anotá el par $(x,y)$ que le corresponde. Conviene hacerlo en una tabla: resultado, cálculo, probabilidad, $x$, $y$.
\item \textbf{Armar la tabla de doble entrada}: cada celda $(x,y)$ es la suma de las probabilidades de los resultados con ese par. Los pares imposibles valen 0, y conviene explicar por qué.
\item \textbf{Control:} la suma de todas las celdas tiene que dar 1. Si no da, hay un resultado mal calculado o faltante.
\item \textbf{Marginales:} se suman las filas, que dan $p_X(x)$, y las columnas, que dan $p_Y(y)$. Se escriben en los márgenes de la tabla.
\end{pasolista}
\end{pasos}
\begin{modelo}{Parcial 2025, Tema 2, Ej.\ 1 (nodos con regla de parada)}
10 nodos: 5 activos (A), 2 en reposo (R), 3 desconectados (D). Se observan de a uno; se para en el primer no activo o al tercer nodo. $X$ = activos observados; $Y$ = desconectados observados.

\textbf{Paso 2:} hay regla de parada, así que se usa un árbol de secuencias.
\textbf{Paso 3:}
\begin{center}
\begin{tabular}{@{}lllcc@{}}
\toprule
Secuencia & Cálculo & Prob. & $x$ & $y$\\\midrule
R & $2/10$ & 0,2000 & 0 & 0\\
D & $3/10$ & 0,3000 & 0 & 1\\
A R & $\frac5{10}\cdot\frac29$ & 0,1111 & 1 & 0\\
A D & $\frac5{10}\cdot\frac39$ & 0,1667 & 1 & 1\\
A A R & $\frac5{10}\cdot\frac49\cdot\frac28$ & 0,0556 & 2 & 0\\
A A D & $\frac5{10}\cdot\frac49\cdot\frac38$ & 0,0833 & 2 & 1\\
A A A & $\frac5{10}\cdot\frac49\cdot\frac38$ & 0,0833 & 3 & 0\\\bottomrule
\end{tabular}
\qquad
\begin{tabular}{@{}c|cc|c@{}}
\toprule
& $y=0$ & $y=1$ & $p_X$\\\midrule
$x=0$ & 0,2000 & 0,3000 & 0,5000\\
$x=1$ & 0,1111 & 0,1667 & 0,2778\\
$x=2$ & 0,0556 & 0,0833 & 0,1389\\
$x=3$ & 0,0833 & 0 & 0,0833\\\midrule
$p_Y$ & 0,4500 & 0,5500 & 1\\\bottomrule
\end{tabular}
\end{center}
\textbf{Pasos 4 y 5:} $(3,1)$ es imposible (con 3 activos ya no queda lugar). La suma total da 1. \checkmark

\textbf{Otro mecanismo} (Temario 1, Ej.\ 1: 2 cartuchos de una caja con 3 azules, 2 rojos y 1 verde, sin reposición). Por conteo, $p(x,y)=\binom3x\binom2y\binom1{2-x-y}/15$; por ejemplo, $p(1,1)=3\cdot2/15=6/15$, y $p(0,0)=0$ porque harían falta 2 verdes.
\end{modelo}
\begin{ojo}
\begin{itemize}
\item Tratar una regla de parada como si fuera ``$n$ fijo'' (hipergeométrica con $n=3$). Hay secuencias que terminan antes.
\item Olvidar las celdas imposibles o no explicar por qué valen 0.
\item No controlar que la suma dé 1.
\end{itemize}
\end{ojo}

% ---------------------------------------------------------------------
\subsection{Probabilidades de eventos y probabilidades condicionales}\label{p:condicional}
\begin{cuando}
``Se sabe que\ldots, ¿cuál es la probabilidad de que\ldots?'', ``si $X$ es mayor que $Y$\ldots'', ``escribir la interpretación simbólica en términos de $X$ e $Y$'', $E(Y\,|\,X=x)$.
\end{cuando}
\begin{pasos}
\begin{pasolista}
\item \textbf{Traducir} cada evento a una condición sobre $(x,y)$. Ojo con los eventos que no son ni $X$ ni $Y$ directamente: hay que expresarlos con ellas. Por ejemplo, ``se registró un nodo desconectado'' es $\{Y=0,\ X\le2\}$.
\item \textbf{Escribir la expresión simbólica}: $P(A\,|\,B)=\dfrac{P(A\cap B)}{P(B)}$, con $B$ = lo que ``se sabe''.
\item \textbf{Marcar en la tabla} las celdas de $B$ (denominador) y, entre ellas, las que además cumplen $A$ (numerador).
\item \textbf{Sumar y dividir.} El resultado tiene que estar entre 0 y 1.
\item \textbf{Para una distribución condicional}: $p_{Y|X}(y|x)=\dfrac{p(x,y)}{p_X(x)}$, que es la fila $x$ dividida por su total. Su esperanza es $\E(Y|X=x)=\sum_y y\,p_{Y|X}(y|x)$.
\end{pasolista}
\end{pasos}
\begin{modelo}{Temario 3, P1c y Parcial 2025 T1, Ej.\ 1c}
\textbf{Tabla dada}, con $P(Y>0\,|\,X>Y)$. \emph{Paso 1:} $B=\{X>Y\}$ son las celdas $(0,-1),(1,-1),(1,0),(2,-1),(2,0),(2,1)$, y $A\cap B=\{(2,1)\}$.
\emph{Paso 4:} $P(B)=0{,}12+0{,}05+0{,}20+0{,}03+0{,}25+0{,}10=0{,}75$, así que $P=\dfrac{0{,}10}{0{,}75}=0{,}1333$.

\textbf{Nodos} (5 A, 3 R, 2 D; $Y$ = en reposo). ``Al menos un activo; probabilidad de que se haya registrado un desconectado.''
\emph{Paso 1:} solo se registra a lo sumo un no activo, así que hubo un D si y solo si se paró por un no activo ($X\le2$) que no era R ($Y=0$). \emph{Paso 2:} $P(Y=0,\,X\le2\,|\,X\ge1)$.
\emph{Paso 4:} $\dfrac{p(1,0)+p(2,0)}{1-p_X(0)}=\dfrac{1/9+1/18}{1/2}=\dfrac13$.
\end{modelo}
\begin{ojo}
\begin{itemize}
\item Dividir por 1 (el total) en lugar de dividir por $P(B)$.
\item Traducir mal el evento, por ejemplo incluir $X=3$ cuando se pregunta por un nodo no activo.
\end{itemize}
\end{ojo}

% ---------------------------------------------------------------------
\subsection{Independencia: cómo demostrarla o refutarla}\label{p:indep}
\begin{cuando}
``¿Son $X$ e $Y$ independientes? Justificar con una demostración o un contraejemplo.''
\end{cuando}
\begin{pasos}
\begin{pasolista}
\item \textbf{Escribir el criterio}: $X$ e $Y$ son independientes $\iff p(x,y)=p_X(x)\,p_Y(y)$ \emph{para todo} par $(x,y)$. En el caso continuo, $f(x,y)=f_X(x)f_Y(y)$.
\item \textbf{Para refutar} (es lo más común), alcanza con \textbf{una} celda donde no se cumpla. Atajos:
\begin{itemize}
\item una celda con $p(x,y)=0$ cuyas marginales son positivas;
\item si ya calculaste $\Cov(X,Y)\neq0$, entonces no son independientes;
\item en el caso continuo, si el soporte no es un rectángulo, no son independientes.
\end{itemize}
\item \textbf{Para demostrar}, hay que verificar \textbf{todas} las celdas, o mostrar que la conjunta se factoriza, $f(x,y)=g(x)h(y)$, sobre un soporte rectangular.
\item \textbf{Escribir la justificación con números}: ``$p(\cdot,\cdot)=\ldots\neq p_X(\cdot)\,p_Y(\cdot)=\ldots$''.
\end{pasolista}
\end{pasos}
\begin{modelo}{Modelo A, Ej.\ 1 y Temario 5, P1}
\textbf{Cajas registradoras:} $p(2,2)=0$, pero $p_X(2)\,p_Y(2)=\frac19\cdot\frac19=\frac1{81}\neq0$. No son independientes.
\textbf{Dos dados} ($X$ = primer dado, $Y$ = máximo): $p(6,1)=0$, pero $p_X(6)\,p_Y(1)=\frac16\cdot\frac1{36}\neq0$. No son independientes.
\end{modelo}
\begin{ojo}
\begin{itemize}
\item Concluir que son independientes porque $\Cov=0$. Es falso: covarianza cero no implica independencia.
\item Verificar una sola celda y concluir que \emph{son} independientes. Hacen falta todas.
\end{itemize}
\end{ojo}

% ---------------------------------------------------------------------
\subsection{Esperanza, covarianza y coeficiente de correlación}\label{p:cov}
\begin{cuando}
``Calcular $\Cov(X,Y)$'', ``coeficiente de correlación'', ``$\E(XY)$'', ``interpretar el resultado''.
\end{cuando}
\begin{pasos}
\begin{pasolista}
\item \textbf{Marginales}: si no están, calculalas (procedimiento \ref{p:conjunta}).
\item \textbf{Medias}: $\E(X)=\sum x\,p_X(x)$ y $\E(Y)=\sum y\,p_Y(y)$.
\item \textbf{Varianzas}: $\E(X^2)=\sum x^2p_X(x)$, luego $\V(X)=\E(X^2)-\E(X)^2$. Lo mismo para $Y$. Si una variable es binomial o Bernoulli, se pueden usar las fórmulas directas: $np$ y $np(1-p)$.
\item \textbf{$\E(XY)=\sum\sum x\,y\,p(x,y)$}: solo aportan las celdas con $x\neq0$ e $y\neq0$. Anotá cada producto.
\item \textbf{Covarianza}: $\Cov(X,Y)=\E(XY)-\E(X)\E(Y)$.
\item \textbf{Correlación}: $\rho=\dfrac{\Cov}{\sigma_X\sigma_Y}$. Controlá que $-1\le\rho\le1$; si no, hay un error de cuentas.
\item \textbf{Interpretar}: el signo indica la dirección (positiva: crecen juntas). Con $|\rho|$ cerca de 1 la relación lineal es fuerte; cerca de 0, débil o nula.
\end{pasolista}
\end{pasos}
\begin{modelo}{Temario 2, P1 (bolillas: $X=1$ si la primera es roja; $Y$ = rojas en 2)}
Tabla: $p(0,0)=\frac{42}{90}$, $p(0,1)=\frac{21}{90}$, $p(1,1)=\frac{21}{90}$, $p(1,2)=\frac6{90}$.
\emph{Pasos 2 y 3:} $\E(X)=0{,}3$, $\V(X)=0{,}21$; \ $\E(Y)=\frac{42+12}{90}=0{,}6$, $\E(Y^2)=\frac{42+24}{90}=0{,}7333$, $\V(Y)=0{,}3733$.
\emph{Paso 4:} $\E(XY)=1\cdot1\cdot\frac{21}{90}+1\cdot2\cdot\frac{6}{90}=0{,}3667$. \
\emph{Paso 5:} $\Cov=0{,}3667-0{,}3\cdot0{,}6=0{,}1867$.
\emph{Paso 6:} $\rho=\dfrac{0{,}1867}{\sqrt{0{,}21\cdot0{,}3733}}=0{,}6667$. \
\emph{Paso 7:} hay una asociación lineal positiva y moderadamente fuerte: si la primera bolilla es roja, hay más rojas en total.
\end{modelo}
\begin{ojo}
\begin{itemize}
\item Calcular $\E(XY)$ como $\E(X)\E(Y)$. Eso solo vale si son independientes.
\item Olvidar restar $\E(X)^2$ al calcular la varianza.
\end{itemize}
\end{ojo}

% ---------------------------------------------------------------------
\subsection{Caso continuo: densidad conjunta}\label{p:continuo}
\begin{cuando}
Se da $f(x,y)$ sobre una región, y se piden independencia, una probabilidad como $P(X\le3Y)$, marginales o la varianza de una función de $X$ e $Y$.
\end{cuando}
\begin{pasos}
\begin{pasolista}
\item \textbf{Dibujar el soporte}, la región donde $f>0$. De ahí salen los límites de integración: los de $y$ pueden depender de $x$, y al revés.
\item \textbf{Verificar} que $\iint f\,dy\,dx=1$. Sirve también como control de los límites.
\item \textbf{Marginales}: $f_X(x)=\int f(x,y)\,dy$, integrando $y$ entre sus límites para $x$ fijo; $f_Y(y)=\int f(x,y)\,dx$, integrando $x$ entre sus límites para $y$ fijo. Al pasar de una a la otra hay que \emph{invertir} la descripción de la región.
\item \textbf{Independencia}: si el soporte no es rectangular, no son independientes. Si lo es, comprobar si $f=f_Xf_Y$.
\item \textbf{Probabilidad de un evento}: dibujar la región del evento intersecada con el soporte y plantear la integral doble con esos límites.
\item \textbf{Momentos}: $\E[g(X,Y)]=\iint g\,f$. Con eso salen $\E(X)$, $\E(X^2)$, $\E(XY)$, luego $\V$ y $\Cov$, y de ahí, por ejemplo, $\V(X-Y)=\V(X)+\V(Y)-2\Cov$.
\end{pasolista}
\end{pasos}
\begin{modelo}{Temario 7, P1: $f(x,y)=2xy$, $0\le x\le2$, $0\le y\le x/2$}
\emph{Paso 2:} $\int_0^2\!\int_0^{x/2}2xy\,dy\,dx=\int_0^2\frac{x^3}{4}dx=1$ \checkmark. \
\emph{Paso 3:} $f_X(x)=\frac{x^3}{4}$ en $[0;2]$; para $y$ fijo, $x$ va de $2y$ a 2, así que $f_Y(y)=\int_{2y}^2 2xy\,dx=4y(1-y^2)$ en $[0;1]$.
\emph{Paso 4:} el soporte es un triángulo, no un rectángulo, así que no son independientes.
\emph{Paso 5:} $P(X\le3Y)$ corresponde a $\frac x3\le y\le\frac x2$: $\int_0^2\!\int_{x/3}^{x/2}2xy\,dy\,dx=\frac5{36}\int_0^2x^3dx=\frac59$.
\emph{Paso 6:} $\V(X)=\frac8{75}$, $\V(Y)=\frac{11}{225}$, $\Cov=\frac8{225}$, así que $\V(X-Y)=\frac{19}{225}=0{,}0844$.
\end{modelo}
\begin{ojo}
\begin{itemize}
\item Poner límites constantes cuando dependen de la otra variable.
\item Calcular $f_Y$ integrando $x$ de 0 a 2: en este ejemplo, $x$ empieza en $2y$.
\end{itemize}
\end{ojo}

% =====================================================================
\newpage
\section{Combinaciones lineales, sumas, TCL y aproximaciones}
% =====================================================================

\subsection{Media, varianza y distribución de una combinación lineal}\label{p:comblin}
\begin{cuando}
``Hallar la varianza de $W=5X-3Y+8$'', ``indique media, varianza y distribución de $w=1{,}3x+y+2$'', ``la holgura $H=B-E$''.
\end{cuando}
\begin{pasos}
\begin{pasolista}
\item \textbf{Escribir la combinación} $W=a_1X_1+\dots+a_kX_k+b$ e identificar cada coeficiente \emph{con su signo}.
\item \textbf{Media} (vale siempre): $\E(W)=a_1\E(X_1)+\dots+a_k\E(X_k)+b$.
\item \textbf{Varianza}: $\V(W)=\sum a_i^2\V(X_i)+2\sum_{i<j}a_ia_j\Cov(X_i,X_j)$. Para dos variables: $\V(aX+bY+c)=a^2\V(X)+b^2\V(Y)+2ab\Cov(X,Y)$.
\begin{itemize}
\item Si son independientes, la covarianza es 0.
\item Si te dan $\rho$ en vez de la covarianza: $\Cov=\rho\,\sigma_X\sigma_Y$.
\item La constante $b$ no cambia la varianza.
\end{itemize}
\item \textbf{Distribución}: si las $X_i$ son \textbf{normales e independientes}, $W$ es normal (\emph{propiedad reproductiva}), con la media y la varianza de los pasos 2 y 3. Si no, solo se conocen la media y la varianza.
\item \textbf{Aclarar los supuestos} que no da el enunciado (por ejemplo, la independencia).
\end{pasolista}
\end{pasos}
\begin{modelo}{Temario 1, Ej.\ 2a y Temario 4, T2}
$\V(X)=2$, $\V(Y)=4$, $\Cov=-2$. \ $\V(5X-3Y+8)=25(2)+9(4)+2(5)(-3)(-2)=50+36+60=146$.

$w=1{,}3x+y+2$, con $x\sim\Nor(\mu;\sigma_1^2)$ e $y\sim\Nor(\mu;\sigma_2^2)$: $\E(w)=2{,}3\mu+2$. Si son independientes, $\V(w)=1{,}69\sigma_1^2+\sigma_2^2$ y $w\sim\Nor(2{,}3\mu+2;\ 1{,}69\sigma_1^2+\sigma_2^2)$.
\end{modelo}
\begin{ojo}
\begin{itemize}
\item No elevar los coeficientes al cuadrado, o restar varianzas: $\V(X-Y)=\V(X)+\V(Y)$ si son independientes.
\item Equivocar el signo del término $2ab\Cov$.
\item Confundir $\V(2X)=4\V(X)$ (una pieza duplicada) con $\V(X_1+X_2)=2\V(X)$ (dos piezas distintas).
\end{itemize}
\end{ojo}

% ---------------------------------------------------------------------
\subsection{Probabilidades de totales y promedios (propiedad reproductiva o TCL)}\label{p:tcl}
\begin{cuando}
``Tiempo total de 150 piezas'', ``consumo de 1 año'', ``peso de una caja de 24'', ``promedio mensual'', ``capacidad del tanque'', ``¿cuánto debe valer $n$ para que\ldots?''.
\end{cuando}
\begin{pasos}
\begin{pasolista}
\item \textbf{Variable individual}: definí $X_i$ (una pieza, un mes, un envase) con su media $\mu$ y su desvío $\sigma$. Si cada unidad es a su vez una suma de partes (tarea 1 + tarea 2, contenido + envase), calculá antes su media y varianza con el procedimiento \ref{p:comblin}.
\item \textbf{Variable de interés}: el \textbf{total} $T=\sum_{i=1}^nX_i$ o el \textbf{promedio} $\Xb=T/n$.
\item \textbf{Supuestos}: las $X_i$ son independientes y tienen la misma distribución. Decilo explícitamente.
\item \textbf{Elegir y nombrar el resultado que se usa}:
\begin{itemize}
\item población \textbf{normal}: el resultado es \emph{exacto}, por la propiedad reproductiva;
\item población no normal o desconocida y $n$ grande ($\ge30$): el resultado es \emph{aproximado}, por el \textbf{TCL}.
\end{itemize}
\item \textbf{Parámetros}: $T\sim\Nor(n\mu;\,n\sigma^2)$ y $\Xb\sim\Nor(\mu;\,\sigma^2/n)$. Para un cambio de escala $Y=aX$: $Y\sim\Nor(a\mu;\,a^2\sigma^2)$.
\item \textbf{Estandarizar y buscar en la tabla}: $z=\dfrac{\text{valor}-\text{media}}{\text{desvío}}$. Luego $P(\cdot<\text{valor})=\Phi(z)$ y $P(\cdot>\text{valor})=1-\Phi(z)$.
\item \textbf{Si piden un percentil} (capacidad, valor que se cumple con probabilidad 0,90): $c=\mu+z_\alpha\sigma$, con $z_\alpha$ tal que $\Phi(z_\alpha)=1-\alpha$.
\item \textbf{Si piden $n$}: plantear la desigualdad de probabilidad, despejar el valor $z$ de la tabla, despejar $n$, \emph{redondear hacia arriba} y verificar con $n$ y con $n-1$.
\end{pasolista}
\end{pasos}
\begin{modelo}{Temario 5, P2 (150 piezas) y Temario 1, Ej.\ 2b (despejar $n$)}
\emph{Paso 1:} $D_i=T_{1i}+T_{2i}$, con $\E(D_i)=12+4=16$ y $\V(D_i)=2^2+1^2=5$. \emph{Paso 2:} $S=\sum_{i=1}^{150}D_i$. \emph{Paso 4:} no se conoce la distribución y $n=150$, así que se usa el TCL.
\emph{Paso 5:} $S\approx\Nor(2400;\,750)$, $\sigma_S=27{,}3861$. \ \emph{Paso 6:} $P(S>2450)=1-\Phi\big(\frac{50}{27{,}3861}\big)=1-\Phi(1{,}8257)=0{,}0339$.

\textbf{Despeje de $n$:} con $T=\sum Y_i$, $\E=0$ y $\V=4n$, se pide $P(|T|>35)\ge0{,}12$. Por el TCL: $2[1-\Phi(\frac{35}{2\sqrt n})]\ge0{,}12\Rightarrow\frac{35}{2\sqrt n}\le z_{0,06}=1{,}5548\Rightarrow n\ge126{,}7\Rightarrow n=127$. Verificación: $n=127$ da 0,1205 \checkmark; $n=126$ da 0,1190.
\end{modelo}
\begin{ojo}
\begin{itemize}
\item Para el total de 24 unidades \emph{distintas} usar $24^2\sigma^2$: lo correcto es $24\sigma^2$.
\item No decir qué teorema se usa: reproductiva si la población es normal, TCL si no lo es.
\item Confundir total con promedio: la varianza es $n\sigma^2$ en un caso y $\sigma^2/n$ en el otro.
\end{itemize}
\end{ojo}

% ---------------------------------------------------------------------
\subsection{Aproximaciones de la binomial (normal o Poisson)}\label{p:aprox}
\begin{cuando}
``De 100 circuitos, ¿probabilidad de que funcionen al menos 96?'', ``si el experimento se repite 90 veces\ldots'', ``dar un resultado aproximado''.
\end{cuando}
\begin{pasos}
\begin{pasolista}
\item \textbf{Identificar la binomial}: $X$ = cantidad de éxitos en $n$ ensayos independientes, cada uno con probabilidad $p$. A veces $p$ hay que calcularlo antes (por ejemplo, la probabilidad de que un circuito funcione).
\item \textbf{Elegir la aproximación}:
\begin{itemize}
\item si $np\ge5$ y $n(1-p)\ge5$: \textbf{normal}, $X\approx\Nor(np;\,np(1-p))$;
\item si el evento es raro ($np<5$ o $n(1-p)<5$): \textbf{Poisson} para la cantidad del evento raro, con $\lambda=n\cdot p_{\text{raro}}$.
\end{itemize}
\item \textbf{Normal con corrección por continuidad}:
\begin{itemize}
\item $P(X\le k)\to$ hasta $k+0{,}5$;
\item $P(X<k)\to$ hasta $k-0{,}5$;
\item $P(X\ge k)\to$ desde $k-0{,}5$;
\item $P(X>k)\to$ desde $k+0{,}5$;
\item $P(X=k)\to$ entre $k-0{,}5$ y $k+0{,}5$.
\end{itemize}
\item \textbf{Poisson}: $P(F\le k)=\sum_{j=0}^{k}e^{-\lambda}\frac{\lambda^j}{j!}$.
\item \textbf{Totales de variables discretas que no son binomiales} (por ejemplo, $\sum Y_i$ con $Y_i\in\{0,1,2\}$): se usa el TCL (procedimiento \ref{p:tcl}) con la misma corrección por continuidad.
\end{pasolista}
\end{pasos}
\begin{modelo}{Modelo B, P1d y Temario 2, P1c}
\textbf{Circuitos:} $P(\text{funciona})=0{,}987712$, así que en $n=100$ se esperan $100(0{,}012288)=1{,}2288<5$ fallas. Normal no; Poisson para las fallas $F\approx\mathrm{Poi}(1{,}2288)$:
$P(W\ge96)=P(F\le4)=0{,}9915$.

\textbf{Bolillas:} $S=\sum_{i=1}^{90}Y_i$, con $\E(S)=54$ y $\sigma_S=5{,}7966$. Corrección: $P(S\ge60)\to$ desde 59,5, y $1-\Phi\big(\frac{59{,}5-54}{5{,}7966}\big)=0{,}1714$.
\end{modelo}
\begin{ojo}
\begin{itemize}
\item Usar la normal cuando $n(1-p)<5$.
\item Olvidar la corrección por continuidad, o aplicarla en el sentido equivocado.
\end{itemize}
\end{ojo}

% =====================================================================
\newpage
\section{Estimación}
% =====================================================================

\subsection{Estimadores: insesgadez, varianza, eficiencia y ECM}\label{p:estim}
\begin{cuando}
``Demostrar que $T$ es un estimador insesgado de\ldots'', ``determinar $b$ para que $Y$ sea insesgado'', ``¿cuál es preferible?'', ``valor que minimiza la varianza'', ``error cuadrático medio''.
\end{cuando}
\begin{pasos}
\begin{pasolista}
\item \textbf{Escribir el estimador} en función de variables cuya media y varianza conocés: $\Xb$ (media $\mu$, varianza $\sigma^2/n$), $X/n$ con $X$ binomial (media $p$, varianza $pq/n$), $X_i^2$ (con $\E(X^2)=\V(X)+\E(X)^2$).
\item \textbf{Calcular $\E(\hat\Theta)$} por linealidad y compararla con el parámetro: el sesgo es $b=\E(\hat\Theta)-\theta$. Si $\E(\hat\Theta)=\theta$ para todo $\theta$, el estimador es \textbf{insesgado}.
\item \textbf{Si hay una constante libre} ($a$, $b$) y piden que sea insesgado: igualar $\E(\hat\Theta)=\theta$ y despejar.
\item \textbf{Calcular $\V(\hat\Theta)$}: con muestras independientes, los coeficientes salen al cuadrado y no hay covarianzas.
\item \textbf{Comparar}: entre insesgados, es mejor (\emph{más eficiente}) el de menor varianza. En general se compara el $\mathrm{ECM}=\V+b^2$. La eficiencia relativa es $\V_1/\V_2$.
\item \textbf{Minimizar la varianza} respecto de una constante: derivar, igualar a 0, despejar y verificar que la derivada segunda sea positiva.
\item \textbf{Preguntas conceptuales}, del tipo ``¿depende de la distribución o de $n$?'': decir qué propiedades se usaron (linealidad de $\E$, independencia) y qué hace falta que exista ($\mu$ y $\sigma^2$ finitos).
\end{pasolista}
\end{pasos}
\begin{modelo}{Modelo A, Ej.\ 5a / Temario 2, P2 / Temario 6, P2 / Modelo B, P2}
$Y=a\Xb_1+(1-a)\Xb_2$ con $n_2=4n_1$. \emph{Paso 2:} $\E(Y)=a\mu+(1-a)\mu=\mu$, insesgado.
\emph{Paso 4:} $\V(Y)=\frac{\sigma^2}{n_1}\big[a^2+\frac{(1-a)^2}{4}\big]$.
\emph{Paso 6:} derivando, $2a-\frac{1-a}{2}=0\Rightarrow a=\frac15$, y $b=\frac45$.

$\hat p_1=\frac12(\frac Xn+\frac Y{3n})$ y $\hat p_2=\frac{X+Y}{4n}$: los dos son insesgados; $\V(\hat p_1)=\frac{pq}{3n}>\V(\hat p_2)=\frac{pq}{4n}$, así que se prefiere $\hat p_2$.

$X\sim U(-a;a)$: $\E(X^2)=\V(X)=\frac{a^2}3$, y entonces $\E\big(\frac3n\sum X_i^2\big)=\frac3n\cdot n\cdot\frac{a^2}{3}=a^2$, insesgado.
\end{modelo}
\begin{ojo}
\begin{itemize}
\item Olvidar la independencia (o los coeficientes al cuadrado) al calcular la varianza.
\item Decir ``es el mejor'' solo porque es insesgado: también hay que comparar la varianza.
\end{itemize}
\end{ojo}

% ---------------------------------------------------------------------
\subsection{Intervalo de confianza para la media}\label{p:icmu}
\begin{cuando}
``Estimar la media / el promedio poblacional con una confianza del\ldots'', ``construir un IC de nivel\ldots''.
\end{cuando}
\begin{pasos}
\begin{pasolista}
\item \textbf{Datos}: $n$, $\xb$, y $\sigma$ (si es conocido) o $s$. Si te dan los datos crudos, calculá $\xb=\frac{\sum x_i}{n}$ y $s^2=\frac{\sum(x_i-\xb)^2}{n-1}$.
\item \textbf{Elegir el caso y justificarlo}:
\begin{itemize}
\item $\sigma$ conocido y (población normal o $n\ge30$): $Z=\frac{\Xb-\mu}{\sigma/\sqrt n}\sim\Nor(0;1)$;
\item $\sigma$ desconocido, población normal y $n$ chico: $T=\frac{\Xb-\mu}{S/\sqrt n}\sim t_{n-1}$;
\item $\sigma$ desconocido y $n\ge30$: $Z$ con $s$ (por el TCL); con $t$ da casi lo mismo.
\end{itemize}
Si el enunciado no dice que la población es normal y $n$ es chico, señalá que falta ese dato.
\item \textbf{Valor crítico}: de la confianza $1-\alpha$ sale $\alpha/2$, y de ahí $z_{\alpha/2}$, o $t_{\alpha/2;\,n-1}$ en la fila $\nu=n-1$.
\item \textbf{Error estándar}: $\sigma/\sqrt n$ o $s/\sqrt n$.
\item \textbf{Error máximo}: $E=(\text{valor crítico})\times(\text{error estándar})$.
\item \textbf{Intervalo}: $[\xb-E;\ \xb+E]$, con unidades.
\item \textbf{Interpretar}: ``con una confianza del $(1-\alpha)100\,\%$, la media poblacional está entre\ldots''. Si te preguntan por un valor, fijate si queda dentro o fuera del intervalo.
\end{pasolista}
\end{pasos}
\begin{modelo}{Modelo B, P3a (frenado) y Temario 6, P3 (DBO)}
\textbf{Frenado:} $n=20$, $\xb=109{,}8$, $s=20$, confianza 90\,\%. \emph{Paso 2:} $\sigma$ desconocido y $n<30$, así que se usa $t_{19}$, y falta el dato de que la población es normal.
\emph{Paso 3:} $t_{0,05;19}=1{,}7291$. \emph{Pasos 4 y 5:} $E=1{,}7291\cdot\frac{20}{\sqrt{20}}=7{,}7329$. \emph{Paso 6:} $[102{,}07;\ 117{,}53]$ m.

\textbf{DBO:} $n=25$, $s^2=0{,}184$, 99\,\%: $t_{0,005;24}=2{,}7969$, $E=2{,}7969\cdot\frac{0{,}4290}{5}=0{,}2400$, así que $[34{,}7600;\ 35{,}2400]$.
\end{modelo}
\begin{ojo}
\begin{itemize}
\item Usar $z$ cuando $\sigma$ es desconocido y $n$ es chico. Es exactamente lo que hace falso el Modelo A, Ej.\ 5c.
\item Buscar $\alpha$ en lugar de $\alpha/2$, o dividir por $n$ en lugar de $\sqrt n$.
\end{itemize}
\end{ojo}

% ---------------------------------------------------------------------
\subsection{Intervalo de confianza para la varianza y el desvío}\label{p:icvar}
\begin{cuando}
``Intervalo de confianza para la varianza / el desvío / la variabilidad''.
\end{cuando}
\begin{pasos}
\begin{pasolista}
\item \textbf{Supuesto obligatorio}: la población es normal. Escribilo.
\item \textbf{Calcular $(n-1)s^2$}. Si tenés los datos, es $\sum(x_i-\xb)^2$.
\item \textbf{Valores críticos} en la fila $\nu=n-1$ de la tabla $\chis$: $\chis_{\alpha/2}$ (el grande) y $\chis_{1-\alpha/2}$ (el chico).
\item \textbf{Intervalo para $\sigma^2$}: $\Big[\dfrac{(n-1)s^2}{\chis_{\alpha/2}}\,;\ \dfrac{(n-1)s^2}{\chis_{1-\alpha/2}}\Big]$. El valor crítico grande va en el extremo de abajo.
\item \textbf{Para $\sigma$}: raíz cuadrada de los dos extremos.
\item \textbf{Interpretar}. Recordá que el intervalo no es simétrico alrededor de $s^2$.
\end{pasolista}
\end{pasos}
\begin{modelo}{Temario 2, P3b y Temario 7, P2a}
$n=16$, $s=12$, 95\,\%. \emph{Paso 2:} $(n-1)s^2=15\cdot144=2160$. \emph{Paso 3:} $\chis_{0,025;15}=27{,}4884$ y $\chis_{0,975;15}=6{,}2621$.
\emph{Paso 4:} $[2160/27{,}4884;\ 2160/6{,}2621]=[78{,}58;\ 344{,}93]$. \emph{Paso 5:} $\sigma\in[8{,}86;\ 18{,}57]$.

Con datos (recipientes, 90\,\%): $\sum(x_i-\xb)^2=0{,}48$, $\chis_{0,05;6}=12{,}5916$, $\chis_{0,95;6}=1{,}6354$, así que $\sigma\in[0{,}1952;\ 0{,}5418]$.
\end{modelo}
\begin{ojo}
\begin{itemize}
\item Invertir los dos $\chis$: en ese caso el extremo de abajo sale mayor que el de arriba.
\item Usar $s$ en lugar de $s^2$, u olvidar sacar la raíz para $\sigma$.
\end{itemize}
\end{ojo}

% ---------------------------------------------------------------------
\subsection{Intervalo de confianza para una proporción (y nivel de un IC dado)}\label{p:icp}
\begin{cuando}
``Estimar la proporción de\ldots con una confianza\ldots'', ``¿con qué nivel de confianza se hizo la estimación $[a;b]$?''.
\end{cuando}
\begin{pasos}
\begin{pasolista}
\item \textbf{Definir la variable}: $X_i=1$ si el individuo tiene la característica, $X_i\sim$ Bernoulli$(p)$. Calculá $\hat p=\frac{\text{éxitos}}{n}$. Ojo con qué proporción piden: por ejemplo, ``a tiempo'' y no ``tarde''.
\item \textbf{Condiciones}: $n\hat p\ge5$ y $n\hat q\ge5$. Así, por el TCL, $\hat P$ es aproximadamente normal.
\item \textbf{Error estándar} (con $\hat p$): $\sqrt{\hat p\hat q/n}$.
\item \textbf{Error máximo e intervalo}: $E=z_{\alpha/2}\sqrt{\hat p\hat q/n}$; el IC es $\hat p\pm E$.
\item \textbf{Si piden el nivel de un IC dado}: $\hat p$ = punto medio, $E$ = semiamplitud, $z=E\big/\sqrt{\hat p\hat q/n}$ y $1-\alpha=2\Phi(z)-1$.
\item \textbf{Interpretar} en el contexto.
\end{pasolista}
\end{pasos}
\begin{modelo}{Parcial 2025 T1, Ej.\ 3d y Temario 4, Ej.\ 2a}
200 pedidos, 40 tarde, así que $\hat p_{\text{a tiempo}}=0{,}80$. $\sqrt{0{,}8\cdot0{,}2/200}=0{,}0283$, \ $E=1{,}96\cdot0{,}0283=0{,}0554$, y el IC es $[0{,}7446;\ 0{,}8554]$.

IC dado $[0{,}5330;\ 0{,}669]$ con $n=120$: $\hat p=0{,}601$, $E=0{,}068$, $\sqrt{0{,}601\cdot0{,}399/120}=0{,}0447$, $z=1{,}5212$, así que $1-\alpha=2\Phi(1{,}5212)-1=0{,}8718\approx87\,\%$.
\end{modelo}
\begin{ojo}
\begin{itemize}
\item En el IC se usa $\hat p$; en la \emph{prueba} se usa $p_0$. No mezclarlos.
\item Estimar la proporción equivocada (la complementaria).
\end{itemize}
\end{ojo}

% ---------------------------------------------------------------------
\subsection{Tamaño de muestra para estimar}\label{p:nest}
\begin{cuando}
``¿Cuántos hay que incluir en la muestra para que el error no supere\ldots?'', ``¿cuántos registros más?'', ``reducir el error a la mitad'', ``amplitud del intervalo de\ldots''.
\end{cuando}
\begin{pasos}
\begin{pasolista}
\item \textbf{Identificar} el parámetro ($\mu$ o $p$), la confianza (de ahí sale $z_{\alpha/2}$) y el error máximo $E$. Si te dan la \emph{amplitud} $A$, entonces $E=A/2$. Si piden ``se desvíe a lo sumo $\varepsilon$ con probabilidad $\ge\gamma$'', es lo mismo con $E=\varepsilon$ y confianza $\gamma$.
\item \textbf{Para $\mu$}: $n=\big(z_{\alpha/2}\,\sigma/E\big)^2$. Si $\sigma$ no se conoce, se usa una estimación previa ($s$ de una muestra piloto o un dato histórico).
\item \textbf{Para $p$}: con una estimación previa, $n=z_{\alpha/2}^2\,\hat p\hat q/E^2$. Sin estimación, o si piden \emph{garantizarlo}, se usa el peor caso $p=0{,}5$: $n=z_{\alpha/2}^2/(4E^2)$.
\item \textbf{Redondear siempre hacia arriba.}
\item \textbf{``¿Cuántos más?''}: $n_{\text{nuevo}}-n_{\text{actual}}$.
\item \textbf{Cambios relativos} (cambia la confianza y/o la amplitud): como $E\propto z/\sqrt n$, se plantea el cociente y $\sigma$ se cancela: $\sqrt{n_2}=\sqrt{n_1}\,\frac{z_2}{z_1}\cdot\frac{E_1}{E_2}$. Para reducir el error a la mitad con la misma confianza, $n$ se multiplica por 4.
\end{pasolista}
\end{pasos}
\begin{modelo}{Modelo A, Ej.\ 2 / Temario 1, Ej.\ 3 / Temario 4, Ej.\ 2b}
\textbf{Teatro} (desvío máximo 0,15 con probabilidad 0,98, sin dato de $p$): $z_{0,01}=2{,}3263$, $n=\frac{2{,}3263^2\cdot0{,}25}{0{,}15^2}=60{,}13$, así que $n=61$.

\textbf{DBO} ($n_1=25$, 90\,\%, amplitud $\varepsilon$; se pide 98\,\% y amplitud $\varepsilon/2$): $\sqrt{n}\ge5\cdot\frac{2{,}3263}{1{,}6449}\cdot2=14{,}1425$, $n\ge200{,}01$, así que $n=201$: 176 registros más.

\textbf{Error a la mitad} con $n=120$: $n'=4\cdot120=480$, es decir, 360 más.
\end{modelo}
\begin{ojo}
\begin{itemize}
\item Redondear hacia abajo, o usar la amplitud como si fuera $E$.
\item Olvidar restar los que ya se tomaron cuando preguntan ``cuántos más''.
\end{itemize}
\end{ojo}

% =====================================================================
\newpage
\section{Pruebas de hipótesis}
% =====================================================================

\subsection{Procedimiento general y planteo de las hipótesis}\label{p:general}
\begin{cuando}
Cualquier ejercicio que diga ``¿hay evidencia significativa\ldots?'', ``plantear la prueba de hipótesis'', ``¿queda confirmada la sospecha?'', ``¿se firma / se adopta / se aprueba?''.
\end{cuando}
\begin{pasos}
\begin{pasolista}
\item \textbf{Parámetro y datos}: decir sobre qué parámetro es la prueba ($\mu$, $\sigma^2$, $p$, $\mu_D$, $\mu_1-\mu_2$) y anotar los datos ($n$, $\xb$, $s$ o $\sigma$, $\hat p$).
\item \textbf{$H_0$}: siempre lleva el ``$=$''. Es el statu quo: lo que se afirma, la norma o lo que vale si no hay cambio.
\item \textbf{$H_1$}: lo que se quiere \emph{demostrar}, lo que se sospecha o lo que lleva a actuar (tomar medidas, adoptar, firmar). Para la dirección, guiate por la tabla de abajo.
\item \textbf{Nivel de significación $\alpha$}: si no lo dan, usá $0{,}05$ y decilo. Si dan ``el riesgo de\ldots indebidamente'', ese riesgo es $\alpha$.
\item \textbf{Estadístico de prueba y su distribución bajo $H_0$}, justificando la elección: normalidad, $\sigma$ conocido o no, TCL.
\item \textbf{Región de rechazo}: con el valor crítico de la tabla. Conviene escribirla también en unidades originales ($\xb<c$, $\hat p<c$).
\item \textbf{Calcular el estadístico} con los datos.
\item \textbf{Decidir y concluir en el contexto}: ``se rechaza / no se rechaza $H_0$''. Luego, ``hay / no hay evidencia significativa de que\ldots''. Si piden el error posible: si se rechazó, es de \textbf{tipo I}; si no se rechazó, de \textbf{tipo II}.
\end{pasolista}
\end{pasos}
\begin{center}
\small
\renewcommand{\arraystretch}{1.25}
\begin{tabular}{@{}>{\raggedright\arraybackslash}p{9cm}l@{}}
\toprule
Palabras del enunciado & $H_1$\\\midrule
``supera'', ``aumentó'', ``es mayor'', ``excede'', ``más de'' & $\theta>\theta_0$\\
``es menor'', ``disminuyó'', ``menos de'', ``se sospecha que es inferior'' & $\theta<\theta_0$\\
``difiere'', ``cambió'', ``verificar si hay que ajustar'' (desvío para cualquier lado) & $\theta\neq\theta_0$\\
``afirma que \emph{al menos} $\theta_0$'' y ``actuará solo si hay pruebas de que no'' & $\theta<\theta_0$\\
``asegura \emph{a lo sumo} $\theta_0$'' y ``¿hay evidencia en contra?'' & $\theta>\theta_0$\\
``aprobará / incorporará / firmará solo si se prueba que cumple'' & $H_1$ = cumple\\
\bottomrule
\end{tabular}
\end{center}
\begin{modelo}{Parcial 2025 T1, Ej.\ 3a (pedidos a tiempo)}
``Afirma que \emph{al menos} el 85\,\% llega a tiempo y solo tomará medidas si encuentra pruebas de que no.''
\emph{Pasos 2 y 3:} $H_0:p=0{,}85$ (o $p\ge0{,}85$) vs.\ $H_1:p<0{,}85$. \emph{Paso 4:} $\alpha=0{,}05$.
\emph{Paso 5:} $Z=\frac{\hat P-0{,}85}{\sqrt{0{,}85\cdot0{,}15/200}}\approx\Nor(0;1)$, por el TCL, con $np_0=170\ge5$ y $nq_0=30\ge5$.
\emph{Paso 6:} se rechaza si $Z<-1{,}6449$, es decir, si $\hat p<0{,}8085$.
\end{modelo}
\begin{ojo}
\begin{itemize}
\item Poner en $H_0$ lo que se quiere demostrar.
\item Decir ``se acepta $H_0$'' o ``se probó $H_0$''. Lo correcto es ``no se rechaza $H_0$''.
\item Elegir una prueba unilateral porque $\xb$ salió de un lado. La dirección la da el enunciado, no los datos.
\end{itemize}
\end{ojo}

% ---------------------------------------------------------------------
\subsection{Prueba para la media}\label{p:pmu}
\begin{pasos}
\begin{pasolista}
\item Plantear $H_0:\mu=\mu_0$ y $H_1$ según el procedimiento \ref{p:general}.
\item \textbf{Elegir el estadístico}:
\begin{itemize}
\item $\sigma$ conocido (normal o $n\ge30$): $Z_0=\dfrac{\xb-\mu_0}{\sigma/\sqrt n}$;
\item $\sigma$ desconocido, población normal y $n$ chico: $T_0=\dfrac{\xb-\mu_0}{s/\sqrt n}\sim t_{n-1}$;
\item $\sigma$ desconocido y $n\ge30$: $Z_0$ con $s$.
\end{itemize}
\item \textbf{Región de rechazo}: $H_1:\neq$ $\Rightarrow|z_0|>z_{\alpha/2}$; \quad $H_1:>$ $\Rightarrow z_0>z_\alpha$; \quad $H_1:<$ $\Rightarrow z_0<-z_\alpha$. Con $t$, lo mismo usando $t_{\cdot;\,n-1}$.
\item \textbf{En unidades de $\xb$}: $c=\mu_0\pm(\text{valor crítico})\cdot\sigma/\sqrt n$.
\item Calcular $z_0$ o $t_0$, decidir y calcular el valor $P$ (procedimiento \ref{p:valorp}).
\end{pasolista}
\end{pasos}
\begin{modelo}{Temario 7, P4 (manzanas) y Temario 2, P3 (satisfacción)}
\textbf{Manzanas:} $H_0:\mu=90{,}4$ vs.\ $H_1:\mu>90{,}4$ (se quiere probar un aumento); $\sigma=6$ conocido, $n=16$, población normal, así que se usa $Z$.
$z_0=\frac{94{,}3-90{,}4}{6/4}=2{,}60>1{,}6449$: se rechaza $H_0$ y se adopta el tratamiento. $P=1-\Phi(2{,}6)=0{,}0047$.

\textbf{Satisfacción:} $\sigma$ desconocido, $n=16$, población normal, así que se usa $t_{15}$. $t_0=\frac{79-70}{12/4}=3{,}0>t_{0,05;15}=1{,}7531$: se rechaza $H_0$, y se alcanzó el objetivo.
\end{modelo}

% ---------------------------------------------------------------------
\subsection{Prueba para la varianza}\label{p:pvar}
\begin{pasos}
\begin{pasolista}
\item \textbf{Supuesto}: la población es normal. Escribilo.
\item \textbf{Hipótesis}: $H_0:\sigma^2=\sigma_0^2$. Para $H_1$: si se quiere probar que la variabilidad es \emph{aceptable} (menor que el límite), $H_1:\sigma^2<\sigma_0^2$; si se quiere probar que es \emph{excesiva}, $H_1:\sigma^2>\sigma_0^2$.
\item \textbf{Estadístico}: $\chis_0=\dfrac{(n-1)s^2}{\sigma_0^2}\sim\chis_{n-1}$ bajo $H_0$.
\item \textbf{Región de rechazo} (fila $\nu=n-1$): $>$: $\chis_0>\chis_{\alpha}$; \quad $<$: $\chis_0<\chis_{1-\alpha}$; \quad $\neq$: $\chis_0<\chis_{1-\alpha/2}$ o $\chis_0>\chis_{\alpha/2}$.
\item Calcular, decidir, interpretar los errores en el contexto y acotar el valor $P$ con la tabla.
\end{pasolista}
\end{pasos}
\begin{modelo}{Temario 1, Ej.\ 4 (arenas)}
Se quiere probar que la arena es apta, o sea $\sigma^2<2$: $H_0:\sigma^2=2$ vs.\ $H_1:\sigma^2<2$. Se toma $\alpha=0{,}05$ y $n=10$.
$\chis_0=\frac{9\cdot1{,}6}{2}=7{,}2$; región: $\chis_0<\chis_{0,95;9}=3{,}3251$. Como $7{,}2>3{,}3251$, \textbf{no se rechaza}: no hay evidencia de que la arena sea apta. $P=P(\chis_9\le7{,}2)=0{,}3837$.
Tipo I: usar una arena inadecuada. Tipo II: descartar una arena buena.
\end{modelo}

% ---------------------------------------------------------------------
\subsection{Prueba para una proporción}\label{p:pp}
\begin{pasos}
\begin{pasolista}
\item \textbf{Variable}: $X_i\sim$ Bernoulli$(p)$; $\hat p=\frac{x}{n}$. Fijate bien qué proporción se prueba.
\item \textbf{Hipótesis}: $H_0:p=p_0$ y $H_1$ según el enunciado.
\item \textbf{Muestra grande} ($np_0\ge5$ y $nq_0\ge5$): $Z_0=\dfrac{\hat p-p_0}{\sqrt{p_0q_0/n}}$. El error estándar se calcula con $p_0$, porque se supone $H_0$ cierta.
\item \textbf{Región de rechazo}: igual que para la media con $Z$. En unidades de $\hat p$: $c=p_0\pm z\sqrt{p_0q_0/n}$.
\item \textbf{Muestra chica}: se trabaja con $X\sim\mathrm{Bin}(n;p_0)$ exacta. La región se elige en la cola o las colas que indique $H_1$, y $\alpha$ es la probabilidad binomial de esa región.
\item Decidir, concluir y calcular el valor $P$.
\end{pasolista}
\end{pasos}
\begin{modelo}{Parcial 2025 T1, Ej.\ 3b y Temario 5, P4b}
$z_0=\frac{0{,}80-0{,}85}{0{,}0252}=-1{,}9803<-1{,}6449$: se rechaza; hay evidencia de que menos del 85\,\% llega a tiempo. $P=\Phi(-1{,}9803)=0{,}0238$. El error posible es de tipo I (medidas innecesarias).

\textbf{Binomial exacta} ($n=25$, $H_1:p\neq0{,}5$): hacen falta las dos colas, así que corresponde $R_1=\{x\le7\vee x\ge18\}$, con $\alpha=P(X\le7)+P(X\ge18)=0{,}0433$.
\end{modelo}
\begin{ojo}
\begin{itemize}
\item Usar $\hat p$ en el denominador de la prueba: va $p_0$.
\item No verificar $np_0\ge5$ y $nq_0\ge5$.
\end{itemize}
\end{ojo}

% ---------------------------------------------------------------------
\subsection{Valor $P$}\label{p:valorp}
\begin{pasos}
\begin{pasolista}
\item \textbf{Calcular el estadístico observado} ($z_0$, $t_0$ o $\chis_0$).
\item \textbf{Área en la dirección de $H_1$}:
\begin{itemize}
\item $H_1:>$: $P=P(Z\ge z_0)=1-\Phi(z_0)$;
\item $H_1:<$: $P=\Phi(z_0)$;
\item $H_1:\neq$: $P=2[1-\Phi(|z_0|)]$.
\end{itemize}
\item \textbf{Con tablas $t$ o $\chis$}: ubicar el valor observado entre dos columnas de la fila $\nu$ y \emph{acotar} el valor $P$ entre sus $\alpha$. En una prueba bilateral, se multiplica por 2.
\item \textbf{Decidir}: si $P\le\alpha$ se rechaza $H_0$; si $P>\alpha$, no se rechaza. Tiene que coincidir con la decisión por la región de rechazo.
\item \textbf{Si dan el valor $P$ y preguntan la conclusión}: compararlo con el $\alpha$ de la prueba. Si no lo dan, con 0,05, y comentar qué pasaría con 0,01 y 0,10.
\end{pasolista}
\end{pasos}
\begin{modelo}{Temario 6, P4c y Modelo A, Ej.\ 4c}
$z_0=\frac{987-1000}{80/\sqrt{200}}=-2{,}2981$, con $H_1:\mu<1000$: $P=\Phi(-2{,}2981)=0{,}0108<0{,}05$, así que se rechaza y se confirma la sospecha.

Si el valor $P$ es $0{,}011$ y la prueba tiene $\alpha=0{,}0161$: como $0{,}011<0{,}0161$, se rechaza $H_0$.
\end{modelo}
\begin{ojo}
\begin{itemize}
\item Decir que el valor $P$ es ``la probabilidad de que $H_0$ sea cierta''. No lo es.
\item Olvidar el factor 2 en las pruebas bilaterales.
\end{itemize}
\end{ojo}

% ---------------------------------------------------------------------
\subsection{$\alpha$, $\beta$ y potencia a partir de una regla de decisión}\label{p:ab}
\begin{cuando}
``Si se rechaza para $\xb\le78{,}2$, ¿cuál es el nivel de significación?'', ``probabilidad de cometer error tipo II si $\mu=\dots$'', ``probabilidad de incorporarlo / adoptarlo si en realidad\ldots'', ``contextuar los errores''.
\end{cuando}
\begin{pasos}
\begin{pasolista}
\item \textbf{Escribir la regla en unidades originales} ($\xb$ o $\hat p$). Si la dan, se usa esa. Si no, se calcula con $\alpha$: $c=\mu_0\pm z_\alpha\,\sigma/\sqrt n$, o $\pm z_{\alpha/2}$ si es bilateral.
\item \textbf{$\alpha$} $=P(\text{caer en la región de rechazo}\,|\,\theta=\theta_0)$, con la distribución centrada en $\theta_0$.
\item \textbf{$\beta(\theta_1)$} $=P(\text{caer en la región de \emph{aceptación}}\,|\,\theta=\theta_1)$: la misma región, pero con la distribución centrada en $\theta_1$. Para proporciones, el desvío se calcula con $p_1$.
\item \textbf{Potencia} $=1-\beta(\theta_1)$, que es la probabilidad de detectar (rechazar $H_0$) cuando el valor verdadero es $\theta_1$.
\item \textbf{Elegir $\theta_1$ con cuidado}: tiene que estar del lado de $H_1$. Por ejemplo, ``difiere en 1,3'' con $H_1:\mu<80$ significa $\mu_1=78{,}7$.
\item \textbf{Contexto}: tipo I es rechazar $H_0$ siendo cierta, y tipo II es no rechazarla siendo falsa. Escribí qué decisión concreta implica cada uno en el problema.
\end{pasolista}
\end{pasos}
\begin{modelo}{Modelo A, Ej.\ 4 / Temario 4, Ej.\ 1c / Temario 5, P5}
$\sigma/\sqrt n=0{,}84$. \emph{Paso 2:} $\alpha=P(\Xb\le78{,}2\,|\,80)=\Phi(-2{,}1429)=0{,}0161$.
\emph{Pasos 3 y 5:} $\mu_1=78{,}7$, así que $\beta=P(\Xb>78{,}2\,|\,78{,}7)=1-\Phi(-0{,}5952)=0{,}7242$.

\textbf{Azufre:} se incorpora si $\xb<5{,}3529$. Con $\mu=5{,}3$: $P=\Phi\big(\frac{5{,}3529-5{,}3}{0{,}0894}\big)=0{,}7228$, que es la potencia.

\textbf{Límites $\mu_0\pm2\sigma/\sqrt n$:} $\alpha=P(|Z|>2)=0{,}0455$. Con $\mu_v=\mu_0+\sigma/\sqrt n$, $\beta=\Phi(1)-\Phi(-3)=0{,}8400$.
\end{modelo}
\begin{ojo}
\begin{itemize}
\item Centrar en $\mu_0$ el cálculo de $\beta$: hay que centrar en $\mu_1$.
\item Elegir $\mu_1$ del lado de $H_0$.
\item Para proporciones, usar $p_0$ en el desvío de $\beta$: va $p_1$.
\end{itemize}
\end{ojo}

% ---------------------------------------------------------------------
\subsection{Curvas OC y de potencia: elegir o leer un gráfico}\label{p:curvas}
\begin{cuando}
``Indicar el gráfico que corresponde a la curva de la probabilidad del error tipo II'', ``¿el siguiente gráfico representa una curva de potencia?'', ``deducir del gráfico $H_0$, $H_1$ y $\alpha$''.
\end{cuando}
\begin{pasos}
\begin{pasolista}
\item \textbf{Ver qué se grafica}: el eje vertical dice si es $\beta(\theta)$ (curva OC) o la potencia $\pi(\theta)=1-\beta(\theta)$.
\item \textbf{Forma según el tipo de prueba}:
\begin{center}
\renewcommand{\arraystretch}{1.2}
\begin{tabular}{@{}lll@{}}
\toprule
Prueba & Curva OC $\beta(\theta)$ & Potencia $\pi(\theta)$\\\midrule
bilateral ($\neq$) & campana, máximo $1-\alpha$ en $\theta_0$ & valle, mínimo $\alpha$ en $\theta_0$\\
derecha ($H_1:\theta>\theta_0$) & decreciente & creciente\\
izquierda ($H_1:\theta<\theta_0$) & creciente & decreciente\\\bottomrule
\end{tabular}
\end{center}
\item \textbf{Valor en $\theta_0$}: la OC vale $1-\alpha$ y la potencia vale $\alpha$. Esto descarta, por ejemplo, una campana con máximo 0,90 si $\alpha=0{,}05$.
\item \textbf{Desempatar entre opciones parecidas}: calcular $\beta$ en uno o dos puntos con el procedimiento \ref{p:ab} y compararlo con el gráfico.
\item \textbf{Leer un gráfico de potencia}: $\theta_0$ es donde la potencia es mínima; $\alpha$ es ese mínimo; la forma da $H_1$; y $P(\text{no rechazar}\,|\,\theta_1)=1-\pi(\theta_1)=\beta$. Si $\theta_1\neq\theta_0$, esa probabilidad es un error de tipo II.
\item \textbf{Justificar} la elección con la forma, el valor en $\theta_0$ y al menos un punto calculado.
\end{pasolista}
\end{pasos}
\begin{modelo}{Parcial 2025 T1/T2, Ej.\ 3c / Temario 1, Ej.\ 5 / Modelo B, P4b}
\textbf{T1} ($H_1:p<0{,}85$, $\alpha=0{,}05$): la OC es creciente, vale 0,95 en 0,85 y $\beta(0{,}79)=0{,}2607$. Es la \textbf{Opción 3}.
\textbf{T2} (bilateral, $\alpha=0{,}05$): campana con máximo 0,95 en $\mu=20$ y $\beta(19{,}8)=\beta(20{,}2)=0{,}1492$. Es la \textbf{Opción 4}; la 2 tiene máximo 0,90.
\textbf{Temario 1} (potencia simétrica con mínimo 0,05 en $p=0{,}5$): $H_0:p=0{,}5$ vs.\ $H_1:p\neq0{,}5$, $\alpha=0{,}05$, y $\beta(0{,}7)=1-\pi(0{,}7)$ es un error de tipo II.
\textbf{Modelo B} ($H_1:\mu>5$): el gráfico es decreciente y vale $\approx0{,}99$ en $\mu=5$, así que es la OC y \textbf{no} la de potencia.
\end{modelo}

% ---------------------------------------------------------------------
\subsection{Tamaño de muestra para una prueba ($\alpha$ y $\beta$ dados)}\label{p:nprueba}
\begin{cuando}
``Determinar $n$ para que, si la media fuese\ldots, la prueba lo detecte con probabilidad\ldots'', ``el proveedor quiere\ldots y el cliente quiere\ldots'', ``¿alcanza con una muestra de 12?''.
\end{cuando}
\begin{pasos}
\begin{pasolista}
\item \textbf{Traducir los requisitos}: ``rechazar a lo sumo el 5\,\% de los buenos'' da $\alpha=0{,}05$ en $\theta_0$; ``detectar con probabilidad 0,80'' o ``rechazar al menos el 80\,\% de los malos'' da $\beta=0{,}20$ en $\theta_1$; ``con probabilidad 0,98 no avalar a los que no cumplen'' da $\alpha=0{,}02$.
\item \textbf{Plantear las dos condiciones sobre el valor crítico $c$} (en el ejemplo, prueba con $H_1:\mu<\mu_0$):
\[
c=\mu_0-z_\alpha\frac{\sigma}{\sqrt n}\quad(\alpha\text{ en }\mu_0),\qquad c=\mu_1+z_\beta\frac{\sigma}{\sqrt n}\quad(\beta\text{ en }\mu_1).
\]
\item \textbf{Igualar y despejar}: $n=\Big[\dfrac{(z_\alpha+z_\beta)\,\sigma}{|\mu_0-\mu_1|}\Big]^2$. En una prueba bilateral se usa $z_{\alpha/2}$ en lugar de $z_\alpha$. Para proporciones: $n=\Big[\dfrac{z_\alpha\sqrt{p_0q_0}+z_\beta\sqrt{p_1q_1}}{p_0-p_1}\Big]^2$.
\item \textbf{Redondear hacia arriba} y \textbf{verificar} $\beta$ con ese $n$.
\item \textbf{Si solo dan $\alpha$ y una regla fija} (``se detiene si $\xb>5{,}5$, con probabilidad de equivocarse 0,01''): despejar $n$ de $\frac{c-\mu_0}{\sigma/\sqrt n}=z_\alpha$.
\item \textbf{Si preguntan ``¿alcanza con $n$?''}: calcular $c$ con ese $n$, luego $\beta$, y comparar con lo pedido.
\item Escribir la \textbf{regla de decisión final}: ``se rechaza si $\xb<c$''.
\end{pasolista}
\end{pasos}
\begin{modelo}{Temario 2, P4 (partidas) / Parcial 2025 T2, Ej.\ 3e / Temario 3, P2}
\textbf{Partidas:} $\mu_0=1250$, $\mu_1=1100$, $\sigma=180$, $\alpha=0{,}05$, $\beta=0{,}10$: $n=\big[\frac{(1{,}6449+1{,}2816)\,180}{150}\big]^2=12{,}33$, así que $n=13$, y la regla es rechazar si $\xb<1167{,}88$. Verificación: $\beta=0{,}087$ \checkmark.

\textbf{Tornillos} (bilateral, $\delta=0{,}1$, $\sigma=0{,}4$, potencia 0,80): $n=\frac{(1{,}96+0{,}8416)^2\,0{,}16}{0{,}01}=125{,}6$, así que $n=126$.

\textbf{Baterías} (¿alcanza $n=12$?): $c=24\,500+2{,}0537\cdot1299{,}04=27\,167{,}9$; $\beta(29\,400)=\Phi(-1{,}7183)=0{,}0429\le0{,}05$. Sí alcanza.
\end{modelo}
\begin{ojo}
\begin{itemize}
\item Usar $z_\alpha$ en una prueba bilateral: va $z_{\alpha/2}$.
\item Confundir $z_\beta$ con $z_{1-\beta}$: para $\beta=0{,}20$, $z_\beta=0{,}8416$.
\item No verificar el resultado.
\end{itemize}
\end{ojo}

% ---------------------------------------------------------------------
\subsection{Relación entre el intervalo de confianza y la prueba}\label{p:dual}
\begin{pasos}
\begin{pasolista}
\item \textbf{Calcular el IC} que corresponde (procedimientos \ref{p:icmu}--\ref{p:icp}).
\item \textbf{Regla de dualidad}: una prueba \emph{bilateral} con nivel $\alpha$ rechaza $H_0:\theta=\theta_0$ si y solo si $\theta_0$ queda \textbf{fuera} del IC del $(1-\alpha)$.
\item \textbf{Si la prueba es unilateral}, la comparación correcta es con la \textbf{cota} del mismo lado (que usa $z_\alpha$): cota superior si $H_1:<$, cota inferior si $H_1:>$.
\item \textbf{Para proporciones} la dualidad es solo aproximada, porque la prueba usa $p_0$ en el error estándar y el IC usa $\hat p$.
\item \textbf{Concluir}: decir si el IC confirma la decisión de la prueba y, si parece que no, explicar por qué (unilateral contra bilateral, o $p_0$ contra $\hat p$).
\end{pasolista}
\end{pasos}
\begin{modelo}{Parcial 2025 T2, Ej.\ 3d y T1, Ej.\ 3d}
\textbf{Tornillos:} el IC 95\,\% es $[20{,}0193;\ 20{,}2807]$ y no contiene a 20; coincide con el rechazo bilateral.
\textbf{Pedidos:} el IC 95\,\% bilateral $[0{,}7446;\ 0{,}8554]$ contiene a 0,85, pero la prueba unilateral rechazó. La cota superior del 95\,\%, $0{,}80+1{,}6449(0{,}0283)=0{,}8465<0{,}85$, sí coincide con el rechazo.
\end{modelo}

% =====================================================================
\newpage
\section{Teoría: verdadero o falso, demostraciones e interpretación}
% =====================================================================

\subsection{Cómo encarar un verdadero o falso}\label{p:vf}
\begin{cuando}
``Analizar si cada proposición es verdadera o falsa. Si es verdadera, demostrarla; si es falsa, dar un contraejemplo o una explicación clara.''
\end{cuando}
\begin{pasos}
\begin{pasolista}
\item \textbf{Leer las palabras clave}: ``para todo'', ``cualquiera sea la distribución'', ``siempre'', ``$<$'' estricto, ``equidistantes'', ``se debe''. Suelen ser las que hacen falsa la afirmación.
\item \textbf{Intentar demostrarla} con definiciones y propiedades: linealidad de $\E$, $\V(aX+b)=a^2\V(X)$, bilinealidad de $\Cov$, independencia.
\item \textbf{Buscar un caso límite} que la rompa: variable degenerada ($\V=0$), distribución asimétrica, $n$ chico, $p=0{,}5$, $\sigma$ desconocido.
\item \textbf{Si es falsa}: dar un \textbf{contraejemplo concreto con números}, o mostrar el cálculo correcto y compararlo con el del enunciado.
\item \textbf{Si es verdadera}: dar la demostración completa, sin saltear pasos, y nombrar las propiedades que usás.
\item \textbf{Cerrar con una frase}: ``Es falsa porque\ldots'' o ``Es verdadera, como se demostró''.
\end{pasolista}
\end{pasos}
\begin{modelo}{V/F de los parciales y modelos}
\begin{itemize}
\item \textbf{$\Cov(h-X,X)<0$ para todo $h$ y cualquier $X$:} \emph{falsa}. $\Cov(h-X,X)=\Cov(h,X)-\Cov(X,X)=-\V(X)\le0$, y si $X$ es constante ($\V=0$) da 0. Contraejemplo: $P(X=c)=1$.
\item \textbf{Los límites del IC de $\sigma^2$ son equidistantes de $s^2$:} \emph{falsa}. Con $n=10$, $s^2=1$ y 95\,\% el IC es $[0{,}4731;\ 3{,}3329]$, y las distancias a $s^2$ son 0,5269 y 2,3329, porque la $\chis$ es asimétrica.
\item \textbf{Siempre hace falta una muestra preliminar para calcular $n$ en una proporción:} \emph{falsa}. Como $p(1-p)\le\frac14$, el valor $n=z^2/(4\varepsilon^2)$ sirve para cualquier $p$.
\item \textbf{Para reducir la longitud del IC a la mitad hay que cuadruplicar $n$:} \emph{verdadera}. Como $L=2z\sigma/\sqrt n$, para $L/2$ hace falta $\sqrt{n'}=2\sqrt n$, o sea $n'=4n$.
\item \textbf{El IC 97\,\% con $n=9$, $\xb=6$, $s=0{,}12$ es $(5{,}9132;\ 6{,}0868)$:} \emph{falsa}. Ese intervalo se calculó con $z=2{,}1701$; corresponde $t_{0,015;8}=2{,}6338$, que da $(5{,}8946;\ 6{,}1054)$.
\item \textbf{$aX_1+(1-a)X_2$ con $n_2=4n_1$ es insesgado y tiene varianza mínima en $a=\frac15$:} \emph{verdadera}. Se demuestra con el procedimiento \ref{p:estim}.
\end{itemize}
\end{modelo}

% ---------------------------------------------------------------------
\subsection{Demostraciones y deducciones típicas (esquemas)}\label{p:demos}
\begin{pasos}
\small
\begin{enumerate}[label=\textbf{(\arabic*)},leftmargin=*,itemsep=4pt]
\item \textbf{$\mathrm{ECM}=\V+b^2$.} Definición $\mathrm{ECM}=\E[(\hat\Theta-\theta)^2]$. Se suma y resta $\E\hat\Theta$, se desarrolla el cuadrado y el término cruzado se anula porque $\E[\hat\Theta-\E\hat\Theta]=0$. Queda $\V(\hat\Theta)+(\E\hat\Theta-\theta)^2$.
\item \textbf{IC para $\sigma^2$.} Pivote: $\frac{(n-1)S^2}{\sigma^2}\sim\chis_{n-1}$ (población normal). Se plantea $P(\chis_{1-\alpha/2}\le\text{pivote}\le\chis_{\alpha/2})=1-\alpha$. Al invertir se dan vuelta las desigualdades; se multiplica por $(n-1)S^2$ y se obtiene el intervalo.
\item \textbf{IC para $\mu$.} Pivote: $\frac{\Xb-\mu}{\sigma/\sqrt n}\sim\Nor(0;1)$. Se plantea $P(-z_{\alpha/2}\le Z\le z_{\alpha/2})=1-\alpha$ y se despeja $\mu$.
\item \textbf{$\E(S^2)=\sigma^2$.} Se usa $\sum(X_i-\Xb)^2=\sum X_i^2-n\Xb^2$, con $\E(X_i^2)=\sigma^2+\mu^2$ y $\E(\Xb^2)=\frac{\sigma^2}n+\mu^2$. El resultado es $(n-1)\sigma^2$, y al dividir por $n-1$ queda $\sigma^2$.
\item \textbf{Tamaño de muestra para una prueba.} Se escriben $\alpha$ y $\beta$ como condiciones sobre el valor crítico $c$, se igualan y se despeja $n$ (procedimiento \ref{p:nprueba}).
\item \textbf{$\beta(p_v)$ en una prueba de proporción unilateral.} Región $\hat p<p_0-z_\alpha\sqrt{p_0q_0/n}$. Entonces $\beta=1-\Phi\Big(\frac{p_0-p_v-z_\alpha\sqrt{p_0q_0/n}}{\sqrt{p_vq_v/n}}\Big)$, y se iguala el argumento a $z_\beta$ para despejar $n$.
\item \textbf{Enunciado del TCL y cuándo $\Xb$ es normal.} Si las $X_i$ son iid con $\sigma^2$ finita, $\frac{\Xb-\mu}{\sigma/\sqrt n}\to\Nor(0;1)$. $\Xb$ es normal \emph{exacta} si la población es normal (por la propiedad reproductiva) y \emph{aproximada} si $n\ge30$ (por el TCL).
\item \textbf{$\Cov=\E(XY)-\mu_X\mu_Y$; $\V(aX+bY)$; $-1\le\rho\le1$.} Se desarrollan los productos y se usa la linealidad de $\E$. Para $\rho$, se usa que $\V(X/\sigma_X\pm Y/\sigma_Y)=2\pm2\rho\ge0$.
\end{enumerate}
\end{pasos}

% ---------------------------------------------------------------------
\subsection{Interpretación: conclusiones, regiones y gráficos de intervalos}\label{p:interp}
\begin{pasos}
\begin{pasolista}
\item \textbf{Redactar la conclusión correcta}:
\begin{itemize}
\item si se rechaza $H_0$: ``hay evidencia significativa (al nivel $\alpha$) de que [$H_1$ en palabras]'';
\item si no se rechaza: ``no hay evidencia significativa de que [$H_1$]''. Nunca ``se probó $H_0$''.
\end{itemize}
\item \textbf{Elegir una región de rechazo entre opciones}: tiene que estar del lado (o de los lados) que indica $H_1$. Si $H_1$ es bilateral, la región tiene que tener dos colas. Después se calcula el $\alpha$ de cada región.
\item \textbf{Gráfico de 100 intervalos}:
\begin{itemize}
\item contá cuántos intervalos \emph{no} contienen al valor verdadero;
\item con confianza $1-\alpha$ se esperan alrededor de $100\alpha$, porque esa cantidad tiene distribución $\mathrm{Bin}(100;\alpha)$;
\item si hay muchos más, el gráfico no corresponde a ese nivel;
\item verificá también que la línea de referencia esté en el valor verdadero que dice el enunciado.
\end{itemize}
\item \textbf{Errores en contexto}: tipo I y tipo II, siempre como decisiones concretas del problema (``detener la producción sin necesidad'', ``aprobar un producto que no cumple'').
\end{pasolista}
\end{pasos}
\begin{modelo}{Modelo A, Ej.\ 3 y Modelo B, P3b}
$H_0$: el método no es efectivo; se rechazó. La frase ``No hay evidencia significativa para concluir que el método no es efectivo'' está \textbf{mal}: corresponde a no rechazar. La correcta es ``Hay evidencia significativa para concluir que el método es efectivo''.

\textbf{100 intervalos:} hay 14 en rojo. Con 95\,\% se esperan alrededor de 5, y $P(F\ge14)=0{,}0005$, así que el gráfico no corresponde a un nivel de 0,95. Además, la línea de referencia está en 109,39 y no en 110.
\end{modelo}

% =====================================================================
\newpage
\titleformat{\section}{\Large\bfseries\color{azul}}{}{0em}{}
\section{Hoja de fórmulas completa}
% =====================================================================
{\small
\begin{hoja}{A. Probabilidad y distribuciones}
\renewcommand{\arraystretch}{1.45}
\begin{tabular}{@{}>{\raggedright\arraybackslash}p{3.0cm}>{\raggedright\arraybackslash}p{5.4cm}ll@{}}
\toprule
Distribución & $p(x)$ o $f(x)$ & $\E(X)$ & $\V(X)$\\\midrule
Bernoulli$(p)$ & $p^x q^{1-x}$, \ $x\in\{0,1\}$ & $p$ & $pq$\\
Binomial$(n;p)$ & $\binom nx p^xq^{n-x}$, \ $x=0,\dots,n$ & $np$ & $npq$\\
Hipergeométrica $(N;M;n)$ & $\binom Mx\binom{N-M}{n-x}\big/\binom Nn$ & $n\frac MN$ & $n\frac MN\big(1-\frac MN\big)\frac{N-n}{N-1}$\\
Poisson$(\lambda)$ & $e^{-\lambda}\lambda^x/x!$, \ $x=0,1,\dots$ & $\lambda$ & $\lambda$\\
Geométrica$(p)$ (ensayos hasta el 1.\textsuperscript{er} éxito) & $q^{x-1}p$, \ $x=1,2,\dots$ & $1/p$ & $q/p^2$\\
Uniforme $(a;b)$ & $\frac1{b-a}$, \ $a<x<b$ & $\frac{a+b}2$ & $\frac{(b-a)^2}{12}$\\
Exponencial$(\lambda)$ & $\lambda e^{-\lambda x}$, \ $x\ge0$ & $1/\lambda$ & $1/\lambda^2$\\
Normal $\Nor(\mu;\sigma^2)$ & $\frac{1}{\sigma\sqrt{2\pi}}e^{-\frac{(x-\mu)^2}{2\sigma^2}}$ & $\mu$ & $\sigma^2$\\
$t_\nu$ \ / \ $\chis_\nu$ & simétrica en 0 / positiva y asimétrica & $0$ / $\nu$ & $\frac{\nu}{\nu-2}$ / $2\nu$\\
\bottomrule
\end{tabular}

\smallskip
$\E(aX+b)=a\E(X)+b$; \ $\V(aX+b)=a^2\V(X)$; \ $\V(X)=\E(X^2)-\E(X)^2$; \ $\E(X^2)=\V(X)+\E(X)^2$.\\
Estandarización: $Z=\frac{X-\mu}{\sigma}\sim\Nor(0;1)$; \ $P(a<X<b)=\Phi\big(\frac{b-\mu}{\sigma}\big)-\Phi\big(\frac{a-\mu}{\sigma}\big)$; \ $\Phi(-z)=1-\Phi(z)$; \ percentil: $x=\mu+z\sigma$.\\
Chebyshev: $P(|X-\mu|\ge k\sigma)\le\frac1{k^2}$, \ o sea, $P(|X-\mu|<\varepsilon)\ge1-\frac{\sigma^2}{\varepsilon^2}$ (para cualquier distribución).
\end{hoja}

\begin{hoja}{B. Variables aleatorias bidimensionales}
\setlength{\columnsep}{16pt}
\begin{multicols}{2}\raggedright
\textbf{Discreto}\\
$p(x,y)=P(X=x,Y=y)\ge0$, \ $\sum\sum p=1$.\\
$p_X(x)=\sum_y p(x,y)$; \ $p_Y(y)=\sum_x p(x,y)$.\\
$p_{Y|X}(y|x)=\frac{p(x,y)}{p_X(x)}$; \ $\E(Y|x)=\sum_y y\,p_{Y|X}(y|x)$.\\
$P(A|B)=\frac{P(A\cap B)}{P(B)}$.\\
Indep.: $p(x,y)=p_X(x)p_Y(y)$ \ $\forall(x,y)$.\\
$\E[h(X,Y)]=\sum\sum h(x,y)\,p(x,y)$.

\columnbreak
\textbf{Continuo}\\
$f(x,y)\ge0$, \ $\iint f=1$; \ $P[(X,Y)\in A]=\iint_A f$.\\
$f_X(x)=\int f(x,y)\,dy$; \ $f_Y(y)=\int f(x,y)\,dx$.\\
$f_{Y|X}(y|x)=\frac{f(x,y)}{f_X(x)}$.\\
Indep.: $f=f_Xf_Y$ (requiere soporte rectangular).\\
$\E[h(X,Y)]=\iint h(x,y)\,f(x,y)\,dx\,dy$.
\end{multicols}
\vspace{-6pt}
\textbf{Covarianza:} $\Cov(X,Y)=\E(XY)-\mu_X\mu_Y$; \ $\Cov(X,X)=\V(X)$; \ $\Cov(aX+b,cY+d)=ac\,\Cov(X,Y)$; \ $\Cov(h-X,X)=-\V(X)$.\\
\textbf{Correlación:} $\rho=\frac{\Cov(X,Y)}{\sigma_X\sigma_Y}\in[-1;1]$; \ $|\rho|=1\iff Y=aX+b$. \ Independencia $\Rightarrow\Cov=0$; la recíproca es falsa.
\end{hoja}

\begin{hoja}{C. Combinaciones lineales, propiedad reproductiva, TCL y aproximaciones}
\raggedright
$W=\sum a_iX_i+b$: \ $\E(W)=\sum a_i\mu_i+b$; \ $\V(W)=\sum a_i^2\sigma_i^2+2\sum_{i<j}a_ia_j\Cov(X_i,X_j)$.\\
$\V(aX+bY+c)=a^2\V(X)+b^2\V(Y)+2ab\Cov(X,Y)$; \ si son independientes: $\V(X\pm Y)=\V(X)+\V(Y)$.\\
iid: $T=\sum X_i$ tiene $\E=n\mu$ y $\V=n\sigma^2$; \ $\Xb$ tiene $\E=\mu$ y $\V=\sigma^2/n$. \ Población finita: $\V(\Xb)=\frac{\sigma^2}n\cdot\frac{N-n}{N-1}$.\\
\textbf{Reproductiva:} si las $X_i$ son normales independientes, $\sum a_iX_i+b\sim\Nor\big(\sum a_i\mu_i+b;\ \sum a_i^2\sigma_i^2\big)$ (exacto).\\
\textbf{TCL:} si son iid con $\sigma^2$ finita y $n\ge30$: $\Xb\approx\Nor(\mu;\sigma^2/n)$ y $\sum X_i\approx\Nor(n\mu;n\sigma^2)$ (aproximado, cualquier población).\\
\textbf{Binomial a normal} ($np\ge5$, $nq\ge5$): $X\approx\Nor(np;npq)$; \ $\hat P\approx\Nor(p;pq/n)$.\\
\textbf{Corrección por continuidad:} $P(X\le k)\to k+0{,}5$; \ $P(X<k)\to k-0{,}5$; \ $P(X\ge k)\to k-0{,}5$; \ $P(X>k)\to k+0{,}5$.\\
\textbf{Binomial a Poisson} (evento raro, $n$ grande): $X\approx\mathrm{Poi}(\lambda=np)$.\\
Sumas de independientes: $\sum\mathrm{Bin}(n_i;p)=\mathrm{Bin}(\sum n_i;p)$; \ $\sum\mathrm{Poi}(\lambda_i)=\mathrm{Poi}(\sum\lambda_i)$; \ $\sum\chis_{\nu_i}=\chis_{\sum\nu_i}$.
\end{hoja}

\begin{hoja}{D. Muestreo y distribuciones muestrales}
\raggedright
$\xb=\frac1n\sum x_i$; \ $s^2=\frac{\sum(x_i-\xb)^2}{n-1}=\frac{\sum x_i^2-n\xb^2}{n-1}$; \ $\hat p=\frac{x}{n}$.\\[3pt]
\renewcommand{\arraystretch}{1.7}
\begin{tabular}{@{}lll@{}}
\toprule
Estadístico & Distribución & Condición\\\midrule
$Z=\dfrac{\Xb-\mu}{\sigma/\sqrt n}$ & $\Nor(0;1)$ & población normal (exacta) o $n\ge30$ (TCL)\\
$T=\dfrac{\Xb-\mu}{S/\sqrt n}$ & $t_{n-1}$ & población normal, $\sigma$ desconocido\\
$\chis=\dfrac{(n-1)S^2}{\sigma^2}$ & $\chis_{n-1}$ & población normal (no es robusto)\\
$Z=\dfrac{\hat P-p}{\sqrt{pq/n}}$ & $\approx\Nor(0;1)$ & $np\ge5$ y $nq\ge5$\\
$\Xb_1-\Xb_2$ & $\Nor\big(\mu_1-\mu_2;\ \frac{\sigma_1^2}{n_1}+\frac{\sigma_2^2}{n_2}\big)$ & muestras independientes\\
\bottomrule
\end{tabular}
\end{hoja}

\begin{hoja}{E. Estimación puntual}
\raggedright
Sesgo $b=\E(\hat\Theta)-\theta$ (insesgado si $b=0$); \ $\mathrm{ECM}=\E[(\hat\Theta-\theta)^2]=\V(\hat\Theta)+b^2$; \ eficiencia relativa $=\V(\hat\Theta_1)/\V(\hat\Theta_2)$.\\
Consistente: $\mathrm{ECM}\to0$ cuando $n\to\infty$. \ Entre insesgados, el mejor es el de menor varianza.\\
Insesgados: $\Xb\to\mu$ (EE $\sigma/\sqrt n$); \ $\hat P\to p$ (EE $\sqrt{pq/n}$); \ $S^2\to\sigma^2$ ($S$ es sesgado para $\sigma$); \ $\Xb_1-\Xb_2$; \ $\hat P_1-\hat P_2$; \ $S_p^2=\frac{(n_1-1)S_1^2+(n_2-1)S_2^2}{n_1+n_2-2}$.\\
Combinación óptima $a\Xb_1+(1-a)\Xb_2$: \ $a=\frac{n_1}{n_1+n_2}$.
\end{hoja}

\begin{hoja}{F. Intervalos de confianza del $(1-\alpha)100\,\%$}
\renewcommand{\arraystretch}{1.55}
\begin{tabular}{@{}>{\raggedright\arraybackslash}p{1.6cm}>{\raggedright\arraybackslash}p{4.6cm}l@{}}
\toprule
Parámetro & Condiciones & Intervalo\\\midrule
$\mu$ & $\sigma$ conocido (normal o $n\ge30$) & $\xb\pm z_{\alpha/2}\,\sigma/\sqrt n$\\
$\mu$ & $\sigma$ desconocido, normal & $\xb\pm t_{\alpha/2;\,n-1}\,s/\sqrt n$\\
$\mu$ & $\sigma$ desconocido, $n\ge30$ & $\xb\pm z_{\alpha/2}\,s/\sqrt n$\\
$\sigma^2$ & normal & $\big[\frac{(n-1)s^2}{\chis_{\alpha/2;\,n-1}}\,;\ \frac{(n-1)s^2}{\chis_{1-\alpha/2;\,n-1}}\big]$ \ ($\sigma$: raíz)\\
$p$ & $n\hat p\ge5$, $n\hat q\ge5$ & $\hat p\pm z_{\alpha/2}\sqrt{\hat p\hat q/n}$\\
$\mu_1-\mu_2$ & $\sigma_i$ conocidos o $n_i\ge30$ & $(\xb_1-\xb_2)\pm z_{\alpha/2}\sqrt{\sigma_1^2/n_1+\sigma_2^2/n_2}$\\
$\mu_1-\mu_2$ & normales, $\sigma_1=\sigma_2$ desconocidos & $(\xb_1-\xb_2)\pm t_{\alpha/2;\,n_1+n_2-2}\,s_p\sqrt{1/n_1+1/n_2}$\\
$\mu_1-\mu_2$ & normales, $\sigma_1\neq\sigma_2$ (Welch) & $(\xb_1-\xb_2)\pm t_{\alpha/2;\,\nu}\sqrt{s_1^2/n_1+s_2^2/n_2}$, \ $\nu=\frac{(s_1^2/n_1+s_2^2/n_2)^2}{\frac{(s_1^2/n_1)^2}{n_1-1}+\frac{(s_2^2/n_2)^2}{n_2-1}}$\\
$\mu_D$ & datos apareados & $\bar d\pm t_{\alpha/2;\,n-1}\,s_d/\sqrt n$\\
$p_1-p_2$ & muestras grandes & $(\hat p_1-\hat p_2)\pm z_{\alpha/2}\sqrt{\hat p_1\hat q_1/n_1+\hat p_2\hat q_2/n_2}$\\
\bottomrule
\end{tabular}

\smallskip
Error máximo $E$ = semiamplitud = (valor crítico)$\times$(error estándar). \ Cotas unilaterales: usar $\alpha$ en lugar de $\alpha/2$.\\
Nivel de un IC dado (proporción): $\hat p$ = punto medio, $z=E/\sqrt{\hat p\hat q/n}$, \ $1-\alpha=2\Phi(z)-1$.
\end{hoja}

\begin{hoja}[unbreakable]{G. Tamaños de muestra (redondear siempre hacia arriba)}
\raggedright
\textbf{Estimar $\mu$:} $n=\big(z_{\alpha/2}\,\sigma/E\big)^2$ \quad (si se da la amplitud $A$: $E=A/2$).\\
\textbf{Estimar $p$:} $n=\frac{z_{\alpha/2}^2\,\hat p\hat q}{E^2}$ \ (con estimación previa) \quad o \quad $n=\frac{z_{\alpha/2}^2}{4E^2}$ \ (sin dato, o para garantizarlo para cualquier $p$).\\
\textbf{Cambios relativos:} $E\propto\frac{z}{\sqrt n}$, así que $\sqrt{n_2}=\sqrt{n_1}\,\frac{z_2}{z_1}\cdot\frac{E_1}{E_2}$. Error a la mitad: $n\times4$. \ ``¿Cuántos más?'': $n_2-n_1$.\\
\textbf{Prueba para $\mu$ ($\alpha$ y $\beta$ dados):} $n=\Big[\frac{(z_\alpha+z_\beta)\,\sigma}{|\mu_0-\mu_1|}\Big]^2$ (unilateral); en una bilateral, $z_{\alpha/2}$ en lugar de $z_\alpha$.\\
\textbf{Prueba para $p$:} $n=\Big[\frac{z_\alpha\sqrt{p_0q_0}+z_\beta\sqrt{p_1q_1}}{|p_0-p_1|}\Big]^2$. \quad \textbf{Solo $\alpha$ y una regla fija $c$:} $\frac{|c-\mu_0|}{\sigma/\sqrt n}=z_\alpha$.
\end{hoja}

\begin{hoja}{H. Pruebas de hipótesis}
\raggedright
$\alpha=P(\text{rechazar }H_0\,|\,H_0\text{ V})$ (tipo I); \ $\beta=P(\text{no rechazar }H_0\,|\,H_0\text{ F})$ (tipo II); \ potencia $=1-\beta$. \ Con $n$ fijo, $\downarrow\alpha\Rightarrow\uparrow\beta$; si $\uparrow n$, bajan los dos.\\[3pt]
\renewcommand{\arraystretch}{1.75}
\begin{tabular}{@{}lll>{\raggedright\arraybackslash}p{2.0cm}@{}}
\toprule
$H_0$ & Condiciones & Estadístico & Distribución\\\midrule
$\mu=\mu_0$ & $\sigma$ conocido (normal o $n\ge30$) & $Z_0=\dfrac{\xb-\mu_0}{\sigma/\sqrt n}$ & $\Nor(0;1)$\\
$\mu=\mu_0$ & normal, $\sigma$ desconocido & $T_0=\dfrac{\xb-\mu_0}{s/\sqrt n}$ & $t_{n-1}$\\
$\sigma^2=\sigma_0^2$ & normal & $\chis_0=\dfrac{(n-1)s^2}{\sigma_0^2}$ & $\chis_{n-1}$\\
$p=p_0$ & $np_0,\,nq_0\ge5$ & $Z_0=\dfrac{\hat p-p_0}{\sqrt{p_0q_0/n}}$ (¡con $p_0$!) & $\Nor(0;1)$\\
$p=p_0$ & $n$ chico & $X$ = cantidad de éxitos & $\mathrm{Bin}(n;p_0)$\\
$\mu_1-\mu_2=\Delta_0$ & $\sigma_i$ conocidos o $n_i\ge30$ & $Z_0=\dfrac{\xb_1-\xb_2-\Delta_0}{\sqrt{\sigma_1^2/n_1+\sigma_2^2/n_2}}$ & $\Nor(0;1)$\\
$\mu_1-\mu_2=\Delta_0$ & normales, $\sigma_1=\sigma_2$ & $T_0=\dfrac{\xb_1-\xb_2-\Delta_0}{s_p\sqrt{1/n_1+1/n_2}}$ & $t_{n_1+n_2-2}$\\
$\mu_D=\Delta_0$ & apareadas & $T_0=\dfrac{\bar d-\Delta_0}{s_d/\sqrt n}$ & $t_{n-1}$\\
\bottomrule
\end{tabular}

\smallskip
\renewcommand{\arraystretch}{1.35}
\begin{tabular}{@{}lllll@{}}
\toprule
$H_1$ & Rechazo ($Z$; $t$ igual) & Rechazo ($\chis$, $\nu=n-1$) & Valor $P$ ($Z$) & $\beta(\mu_1)$, \ $d=\frac{(\mu_1-\mu_0)\sqrt n}{\sigma}$\\\midrule
$\theta\neq\theta_0$ & $|z_0|>z_{\alpha/2}$ & $\chis_0>\chis_{\alpha/2}$ o $\chis_0<\chis_{1-\alpha/2}$ & $2[1-\Phi(|z_0|)]$ & $\Phi(z_{\alpha/2}-d)-\Phi(-z_{\alpha/2}-d)$\\
$\theta>\theta_0$ & $z_0>z_\alpha$ & $\chis_0>\chis_\alpha$ & $1-\Phi(z_0)$ & $\Phi(z_\alpha-d)$\\
$\theta<\theta_0$ & $z_0<-z_\alpha$ & $\chis_0<\chis_{1-\alpha}$ & $\Phi(z_0)$ & $1-\Phi(-z_\alpha-d)$\\
\bottomrule
\end{tabular}

\smallskip
\textbf{Decisión:} $P\le\alpha\Rightarrow$ se rechaza $H_0$. \quad \textbf{Región en unidades originales:} $c=\mu_0\pm z\,\sigma/\sqrt n$ \ ($p$: $c=p_0\pm z\sqrt{p_0q_0/n}$).\\
\textbf{$\beta$ en dos pasos:} (1) región de aceptación en $\xb$ (o $\hat p$) bajo $H_0$; (2) probabilidad de caer adentro con el valor verdadero $\mu_1$ (o $p_1$, con desvío $\sqrt{p_1q_1/n}$).\\
\textbf{Curvas:} la OC $\beta(\theta)$ vale $1-\alpha$ en $\theta_0$ y la potencia vale $\alpha$. En una prueba bilateral la OC tiene forma de campana; con $H_1:>$ es decreciente y con $H_1:<$ es creciente.\\
\textbf{Dualidad:} una prueba bilateral rechaza $H_0:\theta=\theta_0$ si y solo si $\theta_0$ queda fuera del IC del $(1-\alpha)$. Una prueba unilateral se compara con la cota del mismo lado.
\end{hoja}

\begin{hoja}[unbreakable]{I. Valores críticos más usados y lectura de tablas}
\raggedright
\renewcommand{\arraystretch}{1.15}
\begin{tabular}{@{}lccccccccc@{}}
\toprule
$\alpha$ & 0,20 & 0,10 & 0,06 & 0,05 & 0,025 & 0,02 & 0,01 & 0,005 & 0,001\\\midrule
$z_\alpha$ & 0,8416 & 1,2816 & 1,5548 & 1,6449 & 1,9600 & 2,0537 & 2,3263 & 2,5758 & 3,0902\\\bottomrule
\end{tabular}

\smallskip
Confianza bilateral: 90\,\% $\to$ 1,6449; \ 95\,\% $\to$ 1,96; \ 98\,\% $\to$ 2,3263; \ 99\,\% $\to$ 2,5758. \quad
$\Phi(1)=0{,}8413$; \ $\Phi(2)=0{,}9772$; \ $\Phi(3)=0{,}9987$.\\
\textbf{Tabla $\Phi$:} la fila da el entero y el primer decimal, y la columna el segundo decimal. Para hallar $z$ dada la probabilidad, se busca la probabilidad \emph{dentro} de la tabla.\\
\textbf{Tabla $t$:} fila $\nu=n-1$, columna $\alpha$ (área derecha). Si $\nu$ no figura, se usa el menor más cercano. Para el valor $P$ se acota entre dos columnas (bilateral: $\times2$).\\
\textbf{Tabla $\chis$:} fila $\nu$, columna = área derecha. Las columnas 0,995 a 0,90 dan los valores chicos (cola izquierda).
\end{hoja}
}

\subsection*{J. Tabla de la normal estándar: \ $\Phi(z)=P(Z\le z)$}
Para $z<0$: \ $\Phi(-z)=1-\Phi(z)$. \ Ejemplo: $\Phi(1{,}96)=0{,}9750$ (fila 1,9, columna 0,06).
\begin{center}
\footnotesize\setlength{\tabcolsep}{4.2pt}\renewcommand{\arraystretch}{0.95}
\begin{tabular}{@{}c|*{10}{c}@{}}\toprule
$z$ & 0,00 & 0,01 & 0,02 & 0,03 & 0,04 & 0,05 & 0,06 & 0,07 & 0,08 & 0,09\\\midrule
0,0 & 0,5000 & 0,5040 & 0,5080 & 0,5120 & 0,5160 & 0,5199 & 0,5239 & 0,5279 & 0,5319 & 0,5359\\
0,1 & 0,5398 & 0,5438 & 0,5478 & 0,5517 & 0,5557 & 0,5596 & 0,5636 & 0,5675 & 0,5714 & 0,5753\\
0,2 & 0,5793 & 0,5832 & 0,5871 & 0,5910 & 0,5948 & 0,5987 & 0,6026 & 0,6064 & 0,6103 & 0,6141\\
0,3 & 0,6179 & 0,6217 & 0,6255 & 0,6293 & 0,6331 & 0,6368 & 0,6406 & 0,6443 & 0,6480 & 0,6517\\
0,4 & 0,6554 & 0,6591 & 0,6628 & 0,6664 & 0,6700 & 0,6736 & 0,6772 & 0,6808 & 0,6844 & 0,6879\\
\addlinespace[2pt]
0,5 & 0,6915 & 0,6950 & 0,6985 & 0,7019 & 0,7054 & 0,7088 & 0,7123 & 0,7157 & 0,7190 & 0,7224\\
0,6 & 0,7257 & 0,7291 & 0,7324 & 0,7357 & 0,7389 & 0,7422 & 0,7454 & 0,7486 & 0,7517 & 0,7549\\
0,7 & 0,7580 & 0,7611 & 0,7642 & 0,7673 & 0,7704 & 0,7734 & 0,7764 & 0,7794 & 0,7823 & 0,7852\\
0,8 & 0,7881 & 0,7910 & 0,7939 & 0,7967 & 0,7995 & 0,8023 & 0,8051 & 0,8078 & 0,8106 & 0,8133\\
0,9 & 0,8159 & 0,8186 & 0,8212 & 0,8238 & 0,8264 & 0,8289 & 0,8315 & 0,8340 & 0,8365 & 0,8389\\
\addlinespace[2pt]
1,0 & 0,8413 & 0,8438 & 0,8461 & 0,8485 & 0,8508 & 0,8531 & 0,8554 & 0,8577 & 0,8599 & 0,8621\\
1,1 & 0,8643 & 0,8665 & 0,8686 & 0,8708 & 0,8729 & 0,8749 & 0,8770 & 0,8790 & 0,8810 & 0,8830\\
1,2 & 0,8849 & 0,8869 & 0,8888 & 0,8907 & 0,8925 & 0,8944 & 0,8962 & 0,8980 & 0,8997 & 0,9015\\
1,3 & 0,9032 & 0,9049 & 0,9066 & 0,9082 & 0,9099 & 0,9115 & 0,9131 & 0,9147 & 0,9162 & 0,9177\\
1,4 & 0,9192 & 0,9207 & 0,9222 & 0,9236 & 0,9251 & 0,9265 & 0,9279 & 0,9292 & 0,9306 & 0,9319\\
\addlinespace[2pt]
1,5 & 0,9332 & 0,9345 & 0,9357 & 0,9370 & 0,9382 & 0,9394 & 0,9406 & 0,9418 & 0,9429 & 0,9441\\
1,6 & 0,9452 & 0,9463 & 0,9474 & 0,9484 & 0,9495 & 0,9505 & 0,9515 & 0,9525 & 0,9535 & 0,9545\\
1,7 & 0,9554 & 0,9564 & 0,9573 & 0,9582 & 0,9591 & 0,9599 & 0,9608 & 0,9616 & 0,9625 & 0,9633\\
1,8 & 0,9641 & 0,9649 & 0,9656 & 0,9664 & 0,9671 & 0,9678 & 0,9686 & 0,9693 & 0,9699 & 0,9706\\
1,9 & 0,9713 & 0,9719 & 0,9726 & 0,9732 & 0,9738 & 0,9744 & 0,9750 & 0,9756 & 0,9761 & 0,9767\\
\addlinespace[2pt]
2,0 & 0,9772 & 0,9778 & 0,9783 & 0,9788 & 0,9793 & 0,9798 & 0,9803 & 0,9808 & 0,9812 & 0,9817\\
2,1 & 0,9821 & 0,9826 & 0,9830 & 0,9834 & 0,9838 & 0,9842 & 0,9846 & 0,9850 & 0,9854 & 0,9857\\
2,2 & 0,9861 & 0,9864 & 0,9868 & 0,9871 & 0,9875 & 0,9878 & 0,9881 & 0,9884 & 0,9887 & 0,9890\\
2,3 & 0,9893 & 0,9896 & 0,9898 & 0,9901 & 0,9904 & 0,9906 & 0,9909 & 0,9911 & 0,9913 & 0,9916\\
2,4 & 0,9918 & 0,9920 & 0,9922 & 0,9925 & 0,9927 & 0,9929 & 0,9931 & 0,9932 & 0,9934 & 0,9936\\
\addlinespace[2pt]
2,5 & 0,9938 & 0,9940 & 0,9941 & 0,9943 & 0,9945 & 0,9946 & 0,9948 & 0,9949 & 0,9951 & 0,9952\\
2,6 & 0,9953 & 0,9955 & 0,9956 & 0,9957 & 0,9959 & 0,9960 & 0,9961 & 0,9962 & 0,9963 & 0,9964\\
2,7 & 0,9965 & 0,9966 & 0,9967 & 0,9968 & 0,9969 & 0,9970 & 0,9971 & 0,9972 & 0,9973 & 0,9974\\
2,8 & 0,9974 & 0,9975 & 0,9976 & 0,9977 & 0,9977 & 0,9978 & 0,9979 & 0,9979 & 0,9980 & 0,9981\\
2,9 & 0,9981 & 0,9982 & 0,9982 & 0,9983 & 0,9984 & 0,9984 & 0,9985 & 0,9985 & 0,9986 & 0,9986\\
\addlinespace[2pt]
3,0 & 0,9987 & 0,9987 & 0,9987 & 0,9988 & 0,9988 & 0,9989 & 0,9989 & 0,9989 & 0,9990 & 0,9990\\
3,1 & 0,9990 & 0,9991 & 0,9991 & 0,9991 & 0,9992 & 0,9992 & 0,9992 & 0,9992 & 0,9993 & 0,9993\\
3,2 & 0,9993 & 0,9993 & 0,9994 & 0,9994 & 0,9994 & 0,9994 & 0,9994 & 0,9995 & 0,9995 & 0,9995\\
3,3 & 0,9995 & 0,9995 & 0,9995 & 0,9996 & 0,9996 & 0,9996 & 0,9996 & 0,9996 & 0,9996 & 0,9997\\
3,4 & 0,9997 & 0,9997 & 0,9997 & 0,9997 & 0,9997 & 0,9997 & 0,9997 & 0,9997 & 0,9997 & 0,9998\\
\bottomrule\end{tabular}
\end{center}

\newpage
\subsection*{K. Tabla $t$ de Student: \ $t_{\alpha,\nu}$ con $P(T_\nu>t_{\alpha,\nu})=\alpha$}
\begin{center}
\small\setlength{\tabcolsep}{9pt}\renewcommand{\arraystretch}{0.98}
\begin{tabular}{@{}c|*{5}{c}@{}}\toprule
$\nu$ & $0,10$ & $0,05$ & $0,025$ & $0,01$ & $0,005$\\\midrule
1 & 3,078 & 6,314 & 12,706 & 31,821 & 63,657\\
2 & 1,886 & 2,920 & 4,303 & 6,965 & 9,925\\
3 & 1,638 & 2,353 & 3,182 & 4,541 & 5,841\\
4 & 1,533 & 2,132 & 2,776 & 3,747 & 4,604\\
5 & 1,476 & 2,015 & 2,571 & 3,365 & 4,032\\
\addlinespace[2pt]
6 & 1,440 & 1,943 & 2,447 & 3,143 & 3,707\\
7 & 1,415 & 1,895 & 2,365 & 2,998 & 3,499\\
8 & 1,397 & 1,860 & 2,306 & 2,896 & 3,355\\
9 & 1,383 & 1,833 & 2,262 & 2,821 & 3,250\\
10 & 1,372 & 1,812 & 2,228 & 2,764 & 3,169\\
\addlinespace[2pt]
11 & 1,363 & 1,796 & 2,201 & 2,718 & 3,106\\
12 & 1,356 & 1,782 & 2,179 & 2,681 & 3,055\\
13 & 1,350 & 1,771 & 2,160 & 2,650 & 3,012\\
14 & 1,345 & 1,761 & 2,145 & 2,624 & 2,977\\
15 & 1,341 & 1,753 & 2,131 & 2,602 & 2,947\\
\addlinespace[2pt]
16 & 1,337 & 1,746 & 2,120 & 2,583 & 2,921\\
17 & 1,333 & 1,740 & 2,110 & 2,567 & 2,898\\
18 & 1,330 & 1,734 & 2,101 & 2,552 & 2,878\\
19 & 1,328 & 1,729 & 2,093 & 2,539 & 2,861\\
20 & 1,325 & 1,725 & 2,086 & 2,528 & 2,845\\
\addlinespace[2pt]
21 & 1,323 & 1,721 & 2,080 & 2,518 & 2,831\\
22 & 1,321 & 1,717 & 2,074 & 2,508 & 2,819\\
23 & 1,319 & 1,714 & 2,069 & 2,500 & 2,807\\
24 & 1,318 & 1,711 & 2,064 & 2,492 & 2,797\\
25 & 1,316 & 1,708 & 2,060 & 2,485 & 2,787\\
\addlinespace[2pt]
26 & 1,315 & 1,706 & 2,056 & 2,479 & 2,779\\
27 & 1,314 & 1,703 & 2,052 & 2,473 & 2,771\\
28 & 1,313 & 1,701 & 2,048 & 2,467 & 2,763\\
29 & 1,311 & 1,699 & 2,045 & 2,462 & 2,756\\
30 & 1,310 & 1,697 & 2,042 & 2,457 & 2,750\\
\addlinespace[2pt]
40 & 1,303 & 1,684 & 2,021 & 2,423 & 2,704\\
60 & 1,296 & 1,671 & 2,000 & 2,390 & 2,660\\
120 & 1,289 & 1,658 & 1,980 & 2,358 & 2,617\\
$\infty$ & 1,282 & 1,645 & 1,960 & 2,326 & 2,576\\
\bottomrule\end{tabular}
\end{center}

\newpage
\subsection*{L. Tabla ji-cuadrado: \ $\chis_{\alpha,\nu}$ con $P(\chis_\nu>\chis_{\alpha,\nu})=\alpha$}
Las columnas de la izquierda ($\alpha$ grande) dan los valores \emph{chicos}, de la cola izquierda.
\begin{center}
\footnotesize\setlength{\tabcolsep}{3.6pt}
\begin{tabular}{@{}c|*{10}{c}@{}}\toprule
$\nu$ & $0,995$ & $0,99$ & $0,975$ & $0,95$ & $0,90$ & $0,10$ & $0,05$ & $0,025$ & $0,01$ & $0,005$\\\midrule
1 & 0,000 & 0,000 & 0,001 & 0,004 & 0,016 & 2,706 & 3,841 & 5,024 & 6,635 & 7,879\\
2 & 0,010 & 0,020 & 0,051 & 0,103 & 0,211 & 4,605 & 5,991 & 7,378 & 9,210 & 10,60\\
3 & 0,072 & 0,115 & 0,216 & 0,352 & 0,584 & 6,251 & 7,815 & 9,348 & 11,34 & 12,84\\
4 & 0,207 & 0,297 & 0,484 & 0,711 & 1,064 & 7,779 & 9,488 & 11,14 & 13,28 & 14,86\\
5 & 0,412 & 0,554 & 0,831 & 1,145 & 1,610 & 9,236 & 11,07 & 12,83 & 15,09 & 16,75\\
\addlinespace[2pt]
6 & 0,676 & 0,872 & 1,237 & 1,635 & 2,204 & 10,64 & 12,59 & 14,45 & 16,81 & 18,55\\
7 & 0,989 & 1,239 & 1,690 & 2,167 & 2,833 & 12,02 & 14,07 & 16,01 & 18,48 & 20,28\\
8 & 1,344 & 1,646 & 2,180 & 2,733 & 3,490 & 13,36 & 15,51 & 17,53 & 20,09 & 21,95\\
9 & 1,735 & 2,088 & 2,700 & 3,325 & 4,168 & 14,68 & 16,92 & 19,02 & 21,67 & 23,59\\
10 & 2,156 & 2,558 & 3,247 & 3,940 & 4,865 & 15,99 & 18,31 & 20,48 & 23,21 & 25,19\\
\addlinespace[2pt]
11 & 2,603 & 3,053 & 3,816 & 4,575 & 5,578 & 17,28 & 19,68 & 21,92 & 24,72 & 26,76\\
12 & 3,074 & 3,571 & 4,404 & 5,226 & 6,304 & 18,55 & 21,03 & 23,34 & 26,22 & 28,30\\
13 & 3,565 & 4,107 & 5,009 & 5,892 & 7,042 & 19,81 & 22,36 & 24,74 & 27,69 & 29,82\\
14 & 4,075 & 4,660 & 5,629 & 6,571 & 7,790 & 21,06 & 23,68 & 26,12 & 29,14 & 31,32\\
15 & 4,601 & 5,229 & 6,262 & 7,261 & 8,547 & 22,31 & 25,00 & 27,49 & 30,58 & 32,80\\
\addlinespace[2pt]
16 & 5,142 & 5,812 & 6,908 & 7,962 & 9,312 & 23,54 & 26,30 & 28,85 & 32,00 & 34,27\\
17 & 5,697 & 6,408 & 7,564 & 8,672 & 10,09 & 24,77 & 27,59 & 30,19 & 33,41 & 35,72\\
18 & 6,265 & 7,015 & 8,231 & 9,390 & 10,86 & 25,99 & 28,87 & 31,53 & 34,81 & 37,16\\
19 & 6,844 & 7,633 & 8,907 & 10,12 & 11,65 & 27,20 & 30,14 & 32,85 & 36,19 & 38,58\\
20 & 7,434 & 8,260 & 9,591 & 10,85 & 12,44 & 28,41 & 31,41 & 34,17 & 37,57 & 40,00\\
\addlinespace[2pt]
21 & 8,034 & 8,897 & 10,28 & 11,59 & 13,24 & 29,62 & 32,67 & 35,48 & 38,93 & 41,40\\
22 & 8,643 & 9,542 & 10,98 & 12,34 & 14,04 & 30,81 & 33,92 & 36,78 & 40,29 & 42,80\\
23 & 9,260 & 10,20 & 11,69 & 13,09 & 14,85 & 32,01 & 35,17 & 38,08 & 41,64 & 44,18\\
24 & 9,886 & 10,86 & 12,40 & 13,85 & 15,66 & 33,20 & 36,42 & 39,36 & 42,98 & 45,56\\
25 & 10,52 & 11,52 & 13,12 & 14,61 & 16,47 & 34,38 & 37,65 & 40,65 & 44,31 & 46,93\\
\addlinespace[2pt]
26 & 11,16 & 12,20 & 13,84 & 15,38 & 17,29 & 35,56 & 38,89 & 41,92 & 45,64 & 48,29\\
27 & 11,81 & 12,88 & 14,57 & 16,15 & 18,11 & 36,74 & 40,11 & 43,19 & 46,96 & 49,64\\
28 & 12,46 & 13,56 & 15,31 & 16,93 & 18,94 & 37,92 & 41,34 & 44,46 & 48,28 & 50,99\\
29 & 13,12 & 14,26 & 16,05 & 17,71 & 19,77 & 39,09 & 42,56 & 45,72 & 49,59 & 52,34\\
30 & 13,79 & 14,95 & 16,79 & 18,49 & 20,60 & 40,26 & 43,77 & 46,98 & 50,89 & 53,67\\
\addlinespace[2pt]
40 & 20,71 & 22,16 & 24,43 & 26,51 & 29,05 & 51,81 & 55,76 & 59,34 & 63,69 & 66,77\\
50 & 27,99 & 29,71 & 32,36 & 34,76 & 37,69 & 63,17 & 67,50 & 71,42 & 76,15 & 79,49\\
60 & 35,53 & 37,48 & 40,48 & 43,19 & 46,46 & 74,40 & 79,08 & 83,30 & 88,38 & 91,95\\
\bottomrule\end{tabular}
\end{center}


\end{document}
